%%%%%% Document class %%%%%%%%
\documentclass{article} % You can change the font size or other options

%%%%%% Page Layout - Size margins %%%%%%%%
\setlength{\oddsidemargin}{.5cm} 
\setlength{\evensidemargin}{.5cm}
\setlength{\textwidth}{15cm} 
\setlength{\textheight}{20cm}
\setlength{\topmargin}{0cm}
\setlength{\parindent}{0pt}
%\linespread{1.1}

%%%%%% General packages %%%%%%%%
\usepackage{parskip}
\usepackage{float}      % Improved interface for floating objects
\usepackage{graphicx}   % Enhanced support for graphics
\usepackage{hyperref}   % Extensive support for hypertext in LaTeX
\usepackage{latexsym}   % A collection of LaTeX symbols
\usepackage{xcolor}     % Driver-independent color extensions for LaTeX
\usepackage{lipsum}     % General Lorem ipsum text
\usepackage{amsfonts}   % More AMS fonts
\usepackage{amsthm}

%%%%%% Math packages %%%%%%%%
\usepackage{amssymb}    % AMS font symbols
\usepackage{amsmath}    % AMS equation environments
\usepackage{bbm}        % Bitmap version of Computer Modern fonts
\usepackage{eucal}      % Euler script fonts
\usepackage{mathabx}    % More mathematical symbols
\usepackage{mathrsfs}   % Rich set of script (calligraphic) fonts
\usepackage{mathtools}

%%%%%% Table and Array packages %%%%%%%%
\usepackage{multirow}   % Create tabular cells spanning multiple rows
\usepackage{array}      % Extending the array and tabular environments
\usepackage{booktabs}   % Enhances the quality of tables
\usepackage{ulem}       % Package for underlining

%%%%%% Bibliography %%%%%%%%
\usepackage[numbers,sort&compress]{natbib}  % Bibliography management

%%%%%% Utilities (Remove in final document) %%%%%%%%
%\usepackage{showkeys}   % Show label, ref, cite keys
\usepackage{todonotes}  % Adding to-do notes

%%%%%% Enumeration and List Control %%%%%%%%
\usepackage{enumerate}  % Enumerate with redefinable labels
\usepackage{enumitem}   % Control layout of itemize, enumerate, description
\usepackage{thmtools}  %Used for restating lemmas.
%\usepackage{parskip}

%%%%%% User-defined Commands %%%%%%%%
\newcommand{\subscript}[2]{$#1 _ #2$}
\newcommand{\concatenate}[2]{#1.#2}

%%%%%% Theorem Environments %%%%%%%%
\newtheorem{theorem}{Theorem}[section]
\newtheorem{definition}[theorem]{Definition}
\newtheorem{notation}[theorem]{Notation}
\newtheorem{proposition}[theorem]{Proposition}
\newtheorem{lemma}[theorem]{Lemma}
\newtheorem{example}[theorem]{Example}
\newtheorem{assumption}{Assumption}
\newtheorem{corollary}[theorem]{Corollary}
\newtheorem{remark}[theorem]{Remark}

\newtheorem{thm}[section]{Theorem}
\newtheorem{prop}[theorem]{Proposition}
\newtheorem{lem}[theorem]{Lemma}
\newtheorem{ex}[theorem]{Example}
\newtheorem{assum}{Assumption}
\newtheorem{cor}[theorem]{Corollary}
\newtheorem{rem}[theorem]{Remark}

\renewcommand{\theassumption}{\Alph{assumption}}

\makeatletter
\newcommand{\settheoremtag}[1]{% \settheoremtag{<tag>}
\let\oldtheassumption\theassumption% Store \thetheorem
\renewcommand{\theassumption}{#1}% Redefine it to a fixed value
\g@addto@macro\endassumption{% At \end{theorem}, ...
\global\let\theassumption\oldtheassumption}% ...restore \thetheorem
}
\makeatother

%%% LaTeX Header for Mathematical Document %%%

%% Common Symbols %%
\def\wt{\widetilde}
\def\wh{\widehat}
\def\wb{\overline}
\def\ul{\underline}

\newcommand{\1}{\mathbbm{1}} % Indicator function

%% Mathematical Operators %%
\DeclareMathOperator*{\argmin}{arg\,min} 
\DeclareMathOperator*{\argmax}{arg\,max} 
\DeclareMathOperator*{\supess}{ess\,sup} 
\DeclareMathOperator*{\infess}{ess\,inf} 

\DeclareMathOperator{\Tr}{Tr} % Trace of a matrix
\DeclareMathOperator{\Corr}{Corr} % Correlation
\DeclareMathOperator{\Cov}{Cov} % Covariance

%% Custom Commands for Norms %%
\newcommand{\norm}[1]{{\vert\kern-0.25ex\vert #1 \vert\kern-0.25ex\vert}}
\newcommand{\bignorm}[1]{{\big\vert\kern-0.25ex\big\vert #1 \big\vert\kern-0.25ex\big\vert}}
\newcommand{\opnorm}[1]{{\vert\kern-0.25ex\vert\kern-0.25ex\vert #1 \vert\kern-0.25ex\vert\kern-0.25ex\vert}}

%% Probability symbol %%
\newcommand{\dd}{\mathrm{d}}
\newcommand{\angles}[1]{\left\langle #1 \right\rangle}
\newcommand{\DeltaH}{\Delta_h}
\newcommand{\PX}{\mathbb{P}}
\newcommand{\E}{\mathbb{E}}
\newcommand{\EX}{\mathbb{E}}

\DeclareMathOperator{\as}{a.s.}

%% Custom Commands for Limits %%
\newcommand{\limst}{\stackrel{st.}{\Longrightarrow}} % Limit in L^1
\newcommand{\limdist}{\to} % Limit in distribution
\newcommand{\limP}{\stackrel{\mathbb{P}}{\Longrightarrow}} % Limit in probability
\newcommand{\limL}{\stackrel{L^1}{\Longrightarrow}} % Limit in L^1
\newcommand{\limLp}{\stackrel{L^{p}}{\Longrightarrow}} % Limit in L^p
\newcommand{\limas}{\stackrel{\as}{\Longrightarrow}} % Limit almost surely

%% Matrix Commands %%
\newcommand*{\transp}[1]{%
\ensuremath{%
\mskip1mu
\prescript{\smash{\mathrm t}}{}{%
    \mathstrut#1%
}%
}%
}

% --- Macros ---
\newcommand{\lp}{\left\{}
\newcommand{\rp}{\right\}}
\newcommand{\diff}{\mathrm{d}}
\newcommand{\R}{\mathbb{R}}
\newcommand{\M}{\mathbb{M}}
\newcommand{\F}{\mathcal{F}}
\newcommand{\Prob}{\mathbb{P}}
\newcommand{\ind}{\mathbf{1}}
\newcommand{\diff}{\mathrm{d}}

%% Mathbb Definitions %%
\def\bbR{\mathbb{R}}
\def\bbB{\mathbb{B}}
\def\bbC{\mathbb{C}}
\def\bbN{\mathbb{N}}
\def\bbE{\mathbb{E}}
\def\bbZ{\mathbb{Z}}
\def\bbH{\mathbb{H}}
\def\bbX{\mathbb{X}}
\def\bbY{\mathbb{Y}}
\def\bbP{\mathbb{P}}

%% Mathbf Definitions %%
\def\bfa{{\mathbf a}} \def\bfb{{\mathbf b}} \def\bfc{{\mathbf c}}
\def\bfd{{\mathbf d}} \def\bfe{{\mathbf e}} \def\bff{{\mathbf f}}
\def\bfg{{\mathbf g}} \def\bfh{{\mathbf h}} \def\bfi{{\mathbf i}}
\def\bfj{{\mathbf j}} \def\bfk{{\mathbf k}} \def\bfl{{\mathbf l}}
\def\bfm{{\mathbf m}} \def\bfn{{\mathbf n}} \def\bfo{{\mathbf o}}
\def\bfp{{\mathbf p}} \def\bfq{{\mathbf q}} \def\bfr{{\mathbf r}}
\def\bfs{{\mathbf s}} \def\bft{{\mathbf t}} \def\bfu{{\mathbf u}}
\def\bfv{{\mathbf v}} \def\bfw{{\mathbf w}} \def\bfx{{\mathbf x}}
\def\bfy{{\mathbf y}} \def\bfz{{\mathbf z}}
\def\bfB{{\mathbf B}} \def\bfW{{\mathbf W}} \def\bfX{{\mathbf X}}

%% Bold Symbols %%
\newcommand{\bftheta}{{\boldsymbol{\theta}}}
\newcommand{\bfbeta}{{\boldsymbol{\beta}}}
\newcommand{\bfvartheta}{{\boldsymbol{\vartheta}}}
\newcommand{\bfnu}{{\boldsymbol{\nu}}}
\newcommand{\bfeta}{{\boldsymbol{\eta}}}
\newcommand{\bfzeta}{{\boldsymbol{\zeta}}}

%% Cal Symbols %%
\newcommand{\calA}{\mathcal{A}} \newcommand{\calB}{\mathcal{B}}
\newcommand{\calC}{\mathcal{C}} \newcommand{\calD}{\mathcal{D}}
\newcommand{\calE}{\mathcal{E}} \newcommand{\calF}{\mathcal{F}}
\newcommand{\calG}{\mathcal{G}} \newcommand{\calH}{\mathcal{H}}
\newcommand{\calI}{\mathcal{I}} \newcommand{\calJ}{\mathcal{J}}
\newcommand{\calK}{\mathcal{K}} \newcommand{\calL}{\mathcal{L}}
\newcommand{\calM}{\mathcal{M}} \newcommand{\calN}{\mathcal{N}}
\newcommand{\calO}{\mathcal{O}} \newcommand{\calP}{\mathcal{P}}
\newcommand{\calQ}{\mathcal{Q}} \newcommand{\calR}{\mathcal{R}}
\newcommand{\calS}{\mathcal{S}} \newcommand{\calT}{\mathcal{T}}
\newcommand{\calU}{\mathcal{U}} \newcommand{\calV}{\mathcal{V}}
\newcommand{\calW}{\mathcal{W}} \newcommand{\calX}{\mathcal{X}}
\newcommand{\calY}{\mathcal{Y}} \newcommand{\calZ}{\mathcal{Z}}

\setlength{\parindent}{0pt}

\begin{document}

\maketitle
\begin{abstract}
    Conditional
\end{abstract}

\tableofcontents

\section{Introduction}

Point processes form a natural probabilistic framework for modeling random collections of events occurring in time or space. Among them, Poisson point processes play a central role in probability theory, both as canonical models of complete randomness and as fundamental building blocks for more complex stochastic systems. Beyond their direct modeling relevance, Poisson random measures serve as the driving noise in a wide class of stochastic constructions, including Lévy processes, jump diffusions, and interacting particle systems. We refer to \citet{kingman1992poisson, daley2003introduction, last2017lectures} for comprehensive treatments. \textcolor{red}{Need here more references maybe, I asked GPT and he gave these}.

In most applications, however, observed event data depart significantly from the complete randomness of a homogeneous Poisson process. Time-varying activity, clustering, or self-excitation are ubiquitous and require richer models. This has motivated the study of non-homogeneous Poisson processes, and self-exciting point processes such as Hawkes processes in \citet{hawkes1971spectra}. In these models, temporal dependence is encoded through the conditional intensity process $(\lambda_t)_{t \geq 0}$, which may depend on time, exogenous randomness, or the past history of the process. From a probabilistic point of view, such processes are most naturally characterized through their intensities and filtrations, rather than through increment properties specific to the Poisson case.

Because explicit distributions are rarely available beyond the Poisson setting, simulation plays a central role in both the analysis and the practical use of point process models. This is true for statistical inference as well as for the study of perturbed or counterfactual scenarios. A common principle underlying most exact simulation methods is the representation of point processes as functionals of an underlying Poisson random measure. In this framework, a point process with (possibly stochastic) intensity $(\lambda_t)$ is constructed by thinning a Poisson random measure on $\mathbb{R}_+ \times \mathbb{R}_+$, retaining points whose marks fall below the graph of the intensity. This acceptance--rejection procedure, usually referred to as thinning, was introduced for non-homogeneous Poisson processes in \citet{lewis1979simulation} and later extended by \citet{ogata1981lewis, ogata1998space} to accommodate history-dependent intensities, including those arising in Hawkes processes.

From a theoretical perspective, thinning is more than a simulation algorithm. It provides a constructive coupling between a point process and an underlying Poisson random measure. This viewpoint has proved useful in a variety of probabilistic contexts, including the study of stability and ergodicity of self-exciting processes in \citet{bremaud1996stability}, the construction of interacting systems driven by common noise by \citet{bremaud1981point}, and stochastic analysis with jumps in \citet{jacod2012discretization}. \textcolor{red}{Just gave some examples, you probably know more, don't hesitate to add}. In particular, Hawkes processes can be defined as solutions to stochastic equations driven by Poisson random measures, making thinning an intrinsic part of their probabilistic structure.

While the unconditional simulation of point processes is now well understood, a different issue arises in many applications, namely that of conditional or counterfactual simulation. Suppose that a realization of a point process $N$ is observed over a fixed time horizon and interpreted as the outcome of some known dynamics. One may then wish to construct an alternative process $\bar N$ associated with a modified intensity, reflecting an intervention or perturbation of the system, while preserving as much as possible of the original randomness. The objective is not to generate an independent sample from the law of $N$, but rather to construct $\bar N$ on the same probability space and Poisson measure as $N$, so as to allow for meaningful pathwise comparisons.

From a probabilistic standpoint, this problem can be formulated as a coupling problem. Given two intensity processes $(\lambda_t)$ and $(\bar\lambda_t)$, how can one construct two point processes $(N,\bar N)$ driven by a common Poisson random measure, such that $N$ has intensity $\lambda_t$ and $\bar N$ has intensity $\bar\lambda_t$? Classical thinning algorithms only partially address this question, as they are primarily designed for unconditional simulation and do not exploit the information contained in an already observed realization. Yet such conditional constructions are essential whenever the observed trajectory represents the reference path relative to which a perturbation is defined.

This situation arises naturally in a number of fields, including finance, neuroscience, and epidemiology. In financial markets, for instance, one may observe the flow of orders in a limit order book and ask how the system would have evolved had an additional trading strategy been executed. In this setting, the observed order flow is not to be discarded, but rather serves as the reference path relative to which the counterfactual dynamics are defined. Similar questions arise when assessing the effect of intervention policies in epidemic models or the response of neuronal systems to external stimuli. In all these cases, restarting the simulation independently does not address the relevant counterfactual question.

In this paper, we propose a probabilistic framework for the conditional simulation of point processes based on a shared Poisson random measure. Given an observed realization of a point process $N$ with intensity $(\lambda_t)$, we construct a perturbed process $\bar N$ associated with a modified intensity $(\bar\lambda_t)$ by thinning the same underlying Poisson random measure. This construction induces a strong coupling between $N$ and $\bar N$: events compatible with both intensities are preserved, while discrepancies arise  from the imposed intervention and from an independent Poisson measure, that will be explained later. The resulting procedure yields an exact conditional simulation method that naturally extends to stochastic and history-dependent intensities.

We illustrate the usefulness of this framework through an application to market impact in electronic financial markets. Market impact refers to the effect of a trading strategy on asset prices. More precisely, it is defined as the difference between the price trajectory observed when the strategy is executed and the counterfactual price trajectory that would have prevailed had the strategy not been implemented, see \citet{almgren2001optimal}, and \citet{bouchaud2009markets}. In particular, buy (resp. sell) orders tend to exert upward (resp. downward) pressure on prices. When prices are driven by the dynamics of a limit order book, this difference can be traced back to the evolution of queue sizes and order arrivals. Modeling order flow via Poisson or Hawkes processes as in \citet{bacry2015hawkes, jaisson2015market}, and \citet{jusselin2020noarbitrage}, our conditional simulation methodology provides a direct and computationally efficient way to estimate market impact by simulating counterfactual queue dynamics conditional on the observed one.

The remainder of the paper is organized as follows. Section \ref{sec:primer} recalls the necessary background on Poisson random measures and thinning constructions. Section \ref{sec:conditional_measures} introduces the conditional simulation framework and establishes its main properties. Section \ref{sec:applications} applies the methodology to limit order book dynamics and market impact.
\section{A primer on point processes simulation}
\label{sec:primer}
In this section, we will start by defining Poisson random measures and how to simulate a random Poisson measure with given intensity.

We will then describe a constructive approach for the simulation of inhomogeneous Poisson processes using thinning, it corresponds to an adaption of the rejection method for simulating a Poisson process with prescribed intensity.

Finally, we will discuss in this framework what we mean by shocking a point process and how that leads to considering the conditional law of random Poisson measure that will be the object of the study in the next section.

\subsection{Poisson measures and Poisson processes}

We consider in this section a measurable space $(E, \mathcal{E})$ with a $\sigma$-finite measure $\nu$.

\begin{definition}
    A measure $\mu$ on $(E,\mathcal{E})$ is called a counting measure if for all relatively compact $B \in \mathcal{E}$, $\mu(B) \in \mathbb{N}$.
\end{definition}

In particular, a counting measure has finite mass on all compact subspaces.

\begin{definition}
    A Poisson random measure with intensity $\nu$ is a random variable $\pi$ taking values in counting measures on $(E,\mathcal{E})$ such that:
    \begin{itemize}
        \item For all $(B_i)_{i=1}^p \in \mathcal{E}^p$ disjoint, $(\pi(B_i))_{i=1}^p$ are independent.
        \item For all $B \in \mathcal{E}$, the law of $\pi(B)$ is a Poisson distribution with intensity $\nu(B)$.
    \end{itemize}
\end{definition}

\begin{proposition}\textit{Construction of a Poisson measure with finite intensity}, see \cite[Prop. 2.1.6.]{RPG}
    \leavevmode \\
    Let $\Lambda$ be a finite measure on $(E, \mathcal{E})$. Let $N \sim \text{Pois}(\Lambda(E))$ and let $(X_i)_{i \ge 1}$ be a sequence of i.i.d. random elements in $E$ with law $\mathbb{P}(X_i \in \cdot) = \frac{\Lambda(\cdot)}{\Lambda(E)}$, independent of $N$.
    
    Then, the random measure defined by $\Phi = \sum_{i=1}^N \delta_{X_i}$
    is a Poisson random measure with intensity $\Lambda$.
\end{proposition}

\begin{example} \label{ex:simulation_method}
    The main example of interest to us are random Poisson meaures on $\R_+^2$ with intensity the Lebesgue measure.
    
    In every compact rectangle $[0,T] \times [0,M]$, to simulate the points $Z_i$ supporting the Poisson measure, one can either:
    \begin{itemize}
        \item Draw $N$ from $\mathcal{P}(TM)$, and $((T_i, Z_i))_{i=1}^N$ independent from $\mathcal{U}([0,T] \times [0,M])$.
        \item Draw $(T_{i}-T_{i-1}, Z_i)$ from $\mathcal{E}(M) \otimes \mathcal{U}([0,M])$ start from $T_0 = 0$ and stop at the smallest $N$ such that $T_N > T$.
    \end{itemize}
\end{example}

\begin{figure}[H]
    \centering
    \includegraphics[width=0.9\textwidth]{images/scattering_property.png} 
    \caption{Scattering and description of sampling.}
    \label{fig:my_image}
\end{figure}

We consider now a filtered probability space $(\Omega, \mathcal{F}, (\mathcal{F}_t)_{t \ge 0}, \mathbb{P})$.

\begin{definition}[Counting process]
    A counting process is a stochastic process $N = (N_t)_{t \geq 0}$ taking values in $\mathbb{N}$ that is increasing and right-continuous.
\end{definition}

\begin{definition}\textit{Point process with stochastic intensity}
    \leavevmode \\
    Let $N = (N_t)_{t \ge 0}$ be an $(\mathcal{F}_t)$-adapted counting process, $N$ has an $(\mathcal{F}_t)$-stochastic intensity $\lambda = (\lambda_t)_{t \ge 0}$ if:
    \begin{enumerate}
        \item $\lambda$ is a non-negative $(\mathcal{F}_t)$-predictable process such that $\int_0^t \lambda_s \, ds < \infty$ for all $t \ge 0$ almost surely.
        \item The compensated process $M = (M_t)_{t \ge 0}$ defined by $M_t \coloneqq N_t - \int_0^t \lambda_s \, ds$ is an $(\mathcal{F}_t)$-local martingale.
    \end{enumerate}
\end{definition}

\begin{example}
    This definition extends the class of Poisson inhomogeneous process that have deterministic intensity, where the increments $N_t - N_s$ are independent from $\mathcal{F}_s$, to now include as well counting processes with predictable intensities like the following class of processes that will be of interest to us:
    \begin{itemize}
        \item \underline{Hawkes processes:}\\
        The intensity is weighted sum over the previous events, solving $\lambda_t = \mu + \int_0^{t-} \phi(t-s)dN_s$.
        \item \underline{Markovian:} \\
        The intensity is a measurable function $\lambda$ of the current state $\lambda_t = \lambda(N_{t-})$.
    \end{itemize}
\end{example}

In the next section, we will describe an algorithm that provides existence of a large class of Point processes with stochastic intensity as well as a practical way to simulate trajectories of these processes.

\subsection{Thinning algorithm}

We consider in this section a measurable space $(E, \mathcal{E})$ with a $\sigma$-finite measure $\nu$.

\begin{proposition}\textit{Thinned process}, see \cite[Prop. 2.2.6.]{RPG}
    \leavevmode \\
    Let $\Phi = \sum_{k \in \mathbb{N}} \delta_{X_k}$ be a Poisson random measure with intensity $\nu$ on $(E, \mathcal{E})$, $(U_k)_{k \in \mathbb{N}}$ be i.i.d. random variables following $\mathcal{U}([0,1])$, and $p: E \to [0,1]$ be a measurable retention function.
    
    The thinning $\tilde{\Phi}$, defined by
    \[
    \tilde{\Phi} \coloneqq \sum_{k \in \mathbb{N}} 1_{\{U_k \leq p(X_k)\}} \delta_{X_k}
    \]
    is a Poisson random measure with intensity measure $\tilde{\nu}(du) = p(u) \nu(du)$.
\end{proposition}

\begin{example}\textit{Inhomogeneous Poisson process simulation}
    \leavevmode \\
    Consider a bounded measurable function $f: \mathbb{R}_+ \to \mathbb{R}_+$ and with $M = \|f\|_{\infty}$, let $\Phi$ be a Poisson measure on $\mathbb{R}_+ \times [0,M]$ with Lebesgue intensity.\\
    With retention function $p(t,z) = \frac{f(t)}{M}$, the associated thinning $\tilde{\Phi}$ yields a point process:
    \[
    \forall t \geq 0, N_t \coloneqq \tilde{\Phi}([0,t] \times [0,M])
    \]
    $(N_t)_{t \geq 0}$ is a point process with intensity $(f(t))_{t \geq 0}$ with respect to the filtration $\mathcal{F}_t = \sigma(\tilde{\Phi}_{|[0,t]\times [0,M]})$.\\
    Following the description in Example~\ref{ex:simulation_method}, we can simulate this process by drawing times $(T_k)_{k \in \mathbb{N}}$ with independent increments in $\mathcal{E}(M)$ with marks $(Z_k)_{k \in \mathbb{N}}$ i.i.d. following $\mathcal{U}([0,M])$ and at time $T_k$, we record a jump of the process if $Z_k \leq f(T_k)$. This is illustrated in the figure below.
\end{example}

\begin{figure}[H]
    \centering
    \includegraphics[width=0.9\textwidth]{images/thinning_count.png} 
    \caption{Thinning procedure for a counting process.}
    \label{fig:thinning_count}
\end{figure}

The following lemma allows one to extend the construction above to the simulation of point processes with predictable intensity that we build step by step.

\begin{lemma}
\label{next_jump}
\textit{Law of next jump}
    \leavevmode \\
    Let $N$ be a point process with $(\mathcal{F}_t)_{t \geq 0}$-predictable intensity $(\lambda_t)_{t \geq 0}$.\\
    Let $\lambda_t$ be bounded by M in $[0,\tilde{T}]$ with $\tilde{T} = \inf \{s|N_s \neq N_0\}$, and $(T_n)_{n \in \mathbb{N}^*}$ event times of a Poisson measure on $\R_+$ with intensity $\nu(dx)=Mdx$, $(U_n)_{n \in \mathbb{N}^*}$ i.i.d. following $\mathcal{U}([0,1])$.\\
    Then, the law of the first jump of $N$ is the law of $T_S$ where $S = \inf \{ n \in \mathbb{N} | U_n \leq \frac{\lambda_{T_n}}{M} \}$.
\end{lemma}

The lemma is derived from \cite[Theorem 2.8.11.]{MLH}.

\begin{example}\textit{Hawkes Process simulation}
    \leavevmode \\
    The lemma justifies that the following algorithm produces a process $N$ with predictable intensity $\lambda_t = \mu + \int_0^{t-} \phi(t-s) dN_s$ with $\phi$ non-increasing on a given time interval $[0,T]$:

    \begin{center}
    \begin{minipage}{0.9\textwidth}
    \begin{algorithm}[H]
    \label{alg:hawkes}
        \DontPrintSemicolon
        $t \gets 0, \text{times} \gets []$\;
        \While{$t < T$}{
            $\lambda^* \gets \mu + \sum_{s \in \text{times}} \phi(t-s)$\;
            Draw $\tau \sim \mathcal{E}(\lambda^*)$ and $U \sim \mathcal{U}([0,1])$\;
            $t \gets t + \tau$\;
            \If{$t < T$}{
                $\lambda_t \gets \mu + \sum_{s \in \text{times}} \phi(t-s)$\;
                \If{$U \leq \frac{\lambda_t}{\lambda^*}$}{
                    $\text{times} \gets \text{times} \cup \{t\}$\;
                }
            }
        }
        \caption{Hawkes Process Simulation}
    \end{algorithm}
    \end{minipage}
    \end{center}
    The images below describe visually the algorithm.
    \begin{figure}[H]
    \centering
    \includegraphics[width=0.9\textwidth]{images/hawkes_count.png}
    \caption{Thinning procedure for a Hawkes process.}
    \label{fig:hawkes_count}
    \end{figure}
\end{example}

\subsection{Shocking a point process}

In the previous section, we arrived to a representation of a point process $N$ with predictable intensity $(\lambda_t)_{t\geq0}$ as a process $(\tilde{N}, \tilde{\lambda})$ solving a system of equations, with $\pi$ Poisson measure with Lebesgue intensity:
    \[
    \begin{cases} 
        \tilde{N}_t = \int_0^t \int_0^\infty 1_{z \leq \tilde{\lambda}_s} \pi(ds, dz) \\
        \mathcal{L}((\tilde{\lambda}_s,\tilde{N}_s)_{s \leq t}) = \mathcal{L}((\lambda_s,N_s)_{s \leq t})
    \end{cases}
    \]

We are now interested in processes with predictable intensity $\lambda$ of the form:
    \[
    \begin{cases}
    \lambda_t = f(N, t) \\
    N_t = \int_0^t 1_{z \leq \lambda_s} \pi(ds, dz)
    \end{cases}
    \]
    where for $\mathcal{D}$ the space of cadlag functions, we consider a measurable function $f: \mathcal{D} \times \R_+ \to \R_+$ such that $f(A,t) = f(B,t)$ if $A_{|[0,t[} = B_{|[0,t[}$ for $A,B \in \mathcal{D}$.

    When strong solutions exist for a given probability space where $\pi$ exists, we construct shocked processes as follows.

    \begin{proposition}\textit{Shocked process}
        \leavevmode \\
        For a deterministic $\mathbb{Z}$-valued process, the point process $N$ shocked by $M$ is the point process $\bar{N}$ solving the following system:
            \[
            \begin{cases}
            \bar{\lambda}_t = f(\bar{N}, t) \\
            \bar{N}_t = \int_0^t 1_{z \leq \bar{\lambda}_s} \pi(ds, dz) + M_t
            \end{cases}
            \]
    \end{proposition}
\begin{example}\textit{Shocked Hawkess process}
    \leavevmode \\
    We are illustrating this construction on the previous example path in Figure~\ref{fig:hawkes_count}, where we introduce a shock $M_t := 1_{t \geq 1}$ corresponding to a single jump at time 1.

    Now that we have the times and marks from the original simulation, we plot the new intensity curve with the shocks and record a jump of the shocked process whenever the marks are below this new curve. The additional jumps of $\bar{N}$ not seen in $N$ correspond the marks that are sitting above the original intensity curve and below the shocked intensity curve.
\end{example}
    \begin{figure}[H]
    \centering
    \includegraphics[width=0.9\textwidth]{images/shocked_process.png}
    \caption{Shocking procedure of a Hawkes process.}
    \label{fig:shocked_process}
    \end{figure}
    
One can sample paths of the point process and the shocked process for each realization of the Poisson measure.

However, one can only observe one of $N$ or $\bar{N}$ in the applications that we are interested in, this leads us to the question of how can one sample paths from $\bar{N}$ conditionally on $N$ and conversely.

The link between the two processes is made through the Poisson measure $\pi$, this naturally leads us to the question of sampling from $\pi$ conditionally on the path $N$ and solving this problem is enough to access the conditional distribution of $\bar{N}$ given $N$.

\section{Conditional Poisson Measures}
\label{sec:conditional_measures}
In this section, we provide the precise formalism required to characterize the conditional law of a Poisson random measure given its thinned counterpart. We begin by motivating the necessity of a regularity condition on the conditioning domains through a counter-example, followed by a review of the topological frameworks essential for our main results. 

Specifically, we discuss the notion of a random set and the locality requirements necessary to obtain an explicit description of the conditional law.

\subsection{Scattering Property of Random Measures}

The definition of a Poisson random measure $\pi$ implies what we call the \textbf{scattering property}: 
    \[
    \text{For } B \in \mathcal{B}, \, 
    \begin{cases}
        \pi_{|B} \perp \pi_{|B^c}. \\
        \mathcal{L}(\pi) = \mathcal{L}(\pi_{|B} \,\cup \, \pi'_{|B^c}) \text{ with } \pi \perp \pi'.
    \end{cases}
    \]\\
For a process $(N, \lambda)$ solving the following system of equations as in section 2:
    \[
    \begin{cases}
    \lambda_t = f(N, t) \\
    N_t = \int_0^t 1_{z \leq \lambda_s} \pi(ds, dz)
    \end{cases}
    \]

We would like to apply the scattering property to the random set $\{ (z,t) \in \R_+^2 |z \leq \lambda_t \}$ corresponding to the area under the intensity curve used in the simulation process to describe the law of $\pi$ conditionally on the filtration generated by $N$.

However, if now $B$ is a random set taking values in $\mathcal{E}$, we need to understand the conditions on $B$ for $\pi_{|B}$ and $\pi_{|B^c}$ to be independent conditionally on $B$. For such independence to hold, we will see through the next example that we need $B$ to have only a local dependence on the realization of $\pi$.

\begin{example}\textit{Breakdown of independence}
    \leavevmode \\
    Consider a Poisson point process $\pi$ on $[0,1]^2$ with Lebesgue intensity and let $Z_{(1)} > Z_{(2)} > \dots$ denote the ordered $y$-coordinates of the points in $\pi$. We define the random set $A(\pi)$ as the vertical line passing through the point with the second highest $y$-coordinate:
    \[
    A(\pi) \coloneqq \{ (x, y) \in [0,1]^2 \mid y = Z_{(2)} \}.
    \]
    Then, $\pi|_{A(\pi)}$ and $\pi|_{A(\pi)^c}$ are not independent conditionally on $A(\pi)$.
    
    \begin{figure}[H]
        \centering
        \includegraphics[width=0.85\textwidth]{images/independence_breakdown.png}
        \caption{Visualizing the breakdown of independence due to global ordering.}
        \label{fig:independence_breakdown}
    \end{figure}

    The breakdown arises because the information contained in $A(\pi)$ constrains the possible values of points in $A(\pi)^c$.
    
    We observe a shift in the conditional laws:
    \begin{itemize}
        \item The law of $Z_{(2)}$ given only the set $A(\pi)$ is uniform: $\mathcal{L}(Z_{(2)} \mid A(\pi)) = \mathcal{U}([0,1])$.
        \item However, if we also know the location of the first point $Z_{(1)}$, the law is restricted:
        \[
        \mathcal{L}(Z_{(2)} \mid A(\pi), Z_{(1)}) = \mathcal{U}([0, Z_{(1)}]).
        \]
    \end{itemize}
\end{example}

The property on the random set that allows to fix this failure is that $(A(\pi) \subset S)$ is $\pi_{|S}$-measurable.

This property will be formalized in the next subsection under the name of stopping set.

\subsection{Random sets and measures}

\noindent In order to provide a rigorous foundation for our constructions, we must define appropriate topologies on both the space of counting measures—where Poisson measures take their values—and the collection of closed sets of a topological space, which serves as the state space for random sets.

Throughout this section, $\mathbb{G}$ denotes a locally compact second-countable Haussdorf (l.c.s.h.) space, which means that every point admits a compact neighborhood (locally compact), any two points can be separated by open sets (Haussdorf) and the topology admits a countable basis of open sets (second countable).
We denote by $\mathcal{B}(\mathbb{G})$ the Borel $\sigma$-algebra associated with $\mathbb{G}$.

\begin{definition}\textit{Counting measures}
    \leavevmode \\
    The set of counting measures $\mathbb{M}(\mathbb{G})$ consists of all measures $\mu$ on $(\mathbb{G}, \mathcal{B}(\mathbb{G}))$ such that:
    \[
    \mu(B) \in \mathbb{N} \quad \text{for all relatively compact } B \in \mathcal{B}(\mathbb{G}).
    \]
    In particular, counting measures are locally finite. We equip $\mathbb{M}(\mathbb{G})$ with the $\sigma$-algebra $\mathcal{M}(\mathbb{G})$ generated by the evaluation maps:
    \[
    \pi_B : \mu \mapsto \mu(B), \quad \text{for all } B \in \mathcal{B}(\mathbb{G}).
    \]
\end{definition}

\begin{definition}\textit{Collection of closed sets}, see \cite[Definition 5.1.1.]{TCCS}
    \leavevmode \\
    The set $\mathcal{C}(\mathbb{G})$ of all closed subsets of $\mathbb{G}$ is equipped with the \textbf{Fell topology} $\mathcal{T}(\mathbb{G})$. This is the 'hit-and-miss' topology generated by the basis of sets:
    \[
    C^K_{O_1, \dots, O_n} \coloneqq \left\{ F \in C(\mathbb{G}) \mid F \cap K = \emptyset \text{ and } F \cap O_i \neq \emptyset \text{ for } i=1, \dots, n \right\}
    \]
    where $K \subset \mathbb{G}$ is compact and $O_i \subset \mathbb{G}$ are open. 
    
    We denote by $\mathcal{K}(\mathbb{G}) \subset C(\mathbb{G})$ the collection of compact subsets of $\mathbb{G}$, equipped with the subspace topology.
\end{definition}

\begin{remark}
    The Borel $\sigma$-algebra generated by the Fell topology is equal to the $\sigma$-algebra generated by the 'miss' open sets $C_O$ for $O \subset \mathbb{G}$ open.
\end{remark}

\begin{remark}\textit{The Hausdorff and Fell Topologies}, see \cite{TCCS}[Theorem 5.1.10.]
    \leavevmode \\
    The Fell topology $\mathcal{T}(\mathbb{G})$ relates to the Hausdorff metric in the following ways:
    \begin{itemize}
        \item Compact case: If $\mathbb{G}$ is a compact metric space with distance $d$, then $\mathcal{T}(\mathbb{G})$ is the topology induced by the Hausdorff distance $\rho$:
        \[
        \rho(F, F') \coloneqq \max \left( \sup_{x \in F} d(x, F'), \sup_{x' \in F'} d(x', F) \right).
        \]
        
        \item General case: When $\mathbb{G}$ is l.c.s.h. and $(x_i)_{i \in \mathbb{N}}$ a dense sequence in $\mathbb{G}$, then $\mathcal{T}(\mathbb{G})$ is generated by the metric:
        \[
        \rho(F, F') \coloneqq \sum_{i=1}^{\infty} 2^{-i} \min\left( 1, |d(x_i, F) - d(x_i, F')| \right).
        \]
    \end{itemize}
\end{remark}

Throughout the remainder of this section, $(\Omega, \mathcal{F}, \mathbb{P})$ will be a probability space.

\begin{definition}\textit{Random counting measure}
    \leavevmode \\
    A random counting measure on $(\mathbb{G}, \mathcal{B}(\mathbb{G}))$ is a measurable mapping $\pi: (\Omega, \mathcal{F}) \to (\mathbb{M}(\mathbb{G}), \mathcal{M}(\mathbb{G}))$.\\
    A Poisson random measure with intensity $\nu$ is a random counting measure $\pi$ on $(\mathbb{G}, \mathcal{B}(\mathbb{G}))$ such that the joint law of $(\pi(B))_{B \in \mathcal{B}(\mathbb{G})}$ is as in section 2.
\end{definition}

\begin{definition} \textit{Stopping set}
    \leavevmode \\
    Let $\pi$ be a random counting measure on $(\mathbb{G}, \mathcal{B}(\mathbb{G}))$ and $S: \mathbb{M}(\mathbb{G}) \to \mathcal{K}(\mathbb{G})$ a measurable map.\\
    $S(\pi): \Omega \to \mathcal{K}(\mathbb{G})$ defines a random compact set  and is called a stopping set with respect to $\pi$ if:
    \[
    \{S(\pi) \subset K\} \in \sigma(\pi|_K), \quad \forall K \in \mathcal{K}(\mathbb{G}).
    \]
\end{definition}

We can now introduce the main theorem on Poisson point processes.
\begin{theorem}
\textit{Strong Markov property of Poisson point process}, see \cite[Theorem 12.1.3]{RPG}
\label{thm:strongmarkov}
\leavevmode \\
Let $\pi$ be a Poisson random measure on $\mathbb{G}$ and $S(\pi)$ be a stopping set with respect to $\pi$.\\
Then, for all measurable functions $f : \mathbb{M}(\mathbb{G}) \to \mathbb{R}_+$
\[
\mathbb{E}[f(\pi)] = \mathbb{E}[f((\pi|_{S(\pi)}) \cup (\pi'|_{S(\pi)^c}))]
\]
where $\pi'$ is an independent copy of $\pi$.
\end{theorem}

\begin{definition}
    We will denote $\mathcal{D}$ and $\mathbb{L}$ respectively the space of right-continuous with left-limits functions and left-continuous with right-limits functions from $\mathbb{R}_+$ to $\R^d$.

    $\mathcal{D}$ will be equipped with the Kolmogorov $\sigma$-algebra generated by the projections $(f \in \mathcal{D} \to f(t)\in \R^d)_{t \geq 0}$, and $\mathbb{L}$ by the projections $(f \in \mathbb{L} \to f(t)\in \R^d)_{t \geq 0}$.
\end{definition}

We will not delve in Skorokhod topologies as this is not important for our discussion of measurability, the following remark justifies the use of the simpler definition of the $\sigma$-algebra equipping $\mathcal{D}$ and $\mathbb{L}$.

\begin{remark}\textit{Measurability in càglàd functions}, see \cite[Theorem 11.5.2.]{SPL}
    \leavevmode \\
    The Borel $\sigma$-algebra on $\mathbb{L}$ with any of the Skorohod topologies $J_1, J_2, M_1, M_2$ coincides with the Kolomogorov $\sigma$-algebra generated by the projections $(f \in \mathbb{L} \to f(t) \in \R^d)_{t \geq 0}$.\\
    The same statement holds for $\mathcal{D}$.
\end{remark}

\subsection{Main results} \label{subsec:main_results}

In this section, we will restrict our attention to the space $E := [0,T] \times \{ 1,\cdots,d\} \times \R_+$ and $(\Omega, \mathcal{F}, \mathbb{P})$ a probability space.\\
Following the thinning approach to define point processes, we are interested in point processes that admit a pathwise construction. \\
Let $(\lambda, N): \mathbb{M}(E) \to \mathbb{L} \times \mathcal{D}$ with $(\mu \in \mathbb{M}(E) \to (\lambda^\mu, N^\mu) \in \mathbb{L} \times \mathcal{D})$ be a measurable function which there exists a measurable function $f: \mathcal{D} \times \R_+ \to \R_+$ such that
\begin{equation*}
\begin{cases} 
    \lambda^\mu_t = f(j_{t,*}(N^\mu), t) \\
    N^\mu_t = \int_{E} e_k \ind_{z \leq \lambda^{\mu,k}_s} \ind_{s \leq t} \mu(ds, dk, dz)
\end{cases}
\text{ with } 
j_{t,*} \colon \begin{cases} 
    \mathcal{D}' \to \mathcal{D}' \\
    f \mapsto \begin{cases} f(s) & \text{if } s < t \\ 0 & \text{else} \end{cases}
\end{cases}
\text{the extension by zero.}
\end{equation*}

We suppose moreover the \textit{non-explosion} of the counting process:
$$\lambda \text{ is such that } \forall \mu \in \mathbb{M}(E), \sup_{t \in [0,T]} \sup_{k=1}^d \lambda^{\mu,k}_t < \infty.$$

Let $\pi: \Omega \to \mathbb{M}(E)$ be a Poisson random measure with intensity measure $\nu$ given by:
$$\nu(dt,dk,dz) = dt \otimes \eta(dk) \otimes dz \text{ with } \eta({i}) = 1 \text{ for } i \in \{ 1,\cdots,d\}.$$

The process $N^\pi$ defines a point process with predictable intensity $\lambda^\pi$ with respect to the filtration given by $\mathcal{F}^{N^\pi}_t := \sigma(N^\pi_{|[0,t]})$, we denote $\mathcal{F}^{N^\pi}_\infty := \cup_{t \geq 0} \mathcal{F}_t$.

The thinning region for $\mu \in \mathbb{M}(E), A(\mu) := \{ (s,k,z) \in E | z \leq \lambda^{\mu,k}_s \}$ satisfies the following sequence of lemmas leading to the main result.

\begin{restatable}{lemma}{closuremeasurability}
    \label{lem:closuremeasurability}
    The closure $B(\pi)$ of $A(\pi)$ is a compact random set, defining a measurable map from $\Omega$ to $C(E)$.
\end{restatable}

\begin{restatable}{lemma}{stability}
    \label{lem:stability}
    For $\mu, \mu' \in \mathbb{M}(E)$, $A(\mu)$ verifies $A(\mu) = A(\mu_{|B(\mu)} \cup \mu'_{|B(\mu)^c})$.
    In particular, taking the closure yields $B(\mu) = B(\mu_{|B(\mu)} \cup \mu'_{|B(\mu)^c})$.
\end{restatable}

These two lemmas lead to the following result.

\begin{restatable}{lemma}{stoppingset}
\label{lem:stoppingset}
    The random set $B(\pi)$ is a stopping set with respect to $\pi$.
\end{restatable}

\begin{restatable}{theorem}{conditionallaw}
    \label{thm:conditionallaw}
    $\pi_{B(\pi)}$ and $\pi_{B(\pi)^c}$ are independent conditionally on $\mathcal{F}^{N^\pi}_{\infty}$.

    Moreover, the laws can be described as follows:
    \begin{itemize}
        \item $\mathcal{L}(\pi_{|B(\pi)^c}| \mathcal{F}^{N^\pi}_\infty) = \mathcal{P}oisson(\ind_{B(\pi)^c} \nu)$, which is an independent Poisson measure with intensity $\nu$ that is restricted to $B(\pi)^c$.
        \item With $\pi_{|A(\pi)} = \sum_{n \in I} \delta_{(T_n, k_n, Z_n)}$, $\{(T_n, k_n)\}_{n \in I}$ is $\mathcal{F}^{N^\pi}_\infty$-measurable and $\mathcal{L}(Z_n | \mathcal{F}^{N^\pi}_{\infty}) = \mathcal{U}([0, \lambda^{\pi, k_n}_{T_n}])$.
    \end{itemize}
\end{restatable}

The proofs of this section are done in Appendix \ref{appendix:section_three_tree}.

\section{Applications}
\label{sec:applications}
In this section, we illustrate our conditional Poisson process simulation methodology through a concrete real-world application, with a particular focus on financial markets. Throughout this section, we denote by $L^o$ a metaorder executed exclusively through limit orders, and $N^o$ a metaorder executed using market orders. 
\subsection{Limit Order Book: Conditional queue simulation} \label{subsec:lob_sim}
Our object of interest is the queue dynamics of a limit order book, on either the bid or ask side.


% Specifically, the arrival times of market orders are given by the jump times of a Hawkes process $N$ that has a baseline intensity $\mu \geq 0$ and a non-negative self-exciting kernel $\varphi$, so that its intensity is given by
% \begin{equation*}
% \lambda^N_t = \mu + \int_0^{t-} \varphi(t-s)\, dN_s,
% \end{equation*}

We model the aggregated volume available at one side of the order book. Following the queue-reactive framework introduced in \cite{huang2015simulating}, we write
\begin{equation}
\label{eq:queue}
q_t = L_t - C_t - N_t,
\end{equation}
where $L_t$ denotes the cumulative number of limit orders submitted before time $t$, $C_t$ denotes the cumulative number of cancellations, and $N_t$ the cumulative number of market orders.

We assume $L$, $C$, and $N$ to be time-inhomogeneous Poisson processes with intensities
\begin{equation*}
\lambda^{L}_t = \lambda^L\!\left(q_{t-}\right), \qquad \lambda^{C}_t = \lambda^C\!\left(q_{t-}\right), \qquad \lambda^{N}_t = \lambda^N\!\left(q_{t-}\right),
\end{equation*}
for some given functions $\lambda^L$, $\lambda^C$, and $\lambda^N$.

Throughout, we implicitly assume that the Poisson point measures driving the processes $N$, $L$, and $C$ are independent. More precisely, we assume that there exist independent Poisson point measures $\pi^{L}$, $\pi^{C}$, and $\pi^{N}$ such that
\begin{equation}
\label{eq:def:pointmeasure}
\begin{cases}
L_t = \displaystyle \int_0^t \int_0^\infty \mathbf{1}_{\{z \leq \lambda^{L}_s\}}\, \pi^{L}(ds,dz), \\[1ex]
C_t = \displaystyle \int_0^t \int_0^\infty \mathbf{1}_{\{z \leq \lambda^{C}_s\}}\, \pi^{C}(ds,dz), \\[1ex]
N_t = \displaystyle \int_0^t \int_0^\infty \mathbf{1}_{\{z \leq \lambda^N_s\}}\, \pi^{N}(ds,dz), \\
\end{cases}
\end{equation}
    
We now introduce a metaorder into the model. 
Without loss of generality, we assume that it is executed exclusively through limit orders.
The arrival times of this metaorder are modeled by the jump times of a point process $L^o$. We assume it is adapted to the filtration $(\mathcal F_t)_{t \ge 0}$ generated by $(L,C,N)$, in the sense that $L^o_t$ is $\mathcal F_t$-measurable for all $t$, and that its future increments $(L^o_s - L^o_t)_{s \ge t}$ are independent of $\mathcal F_t$.
We further assume that $L^o$ is a non-homogeneous Poisson process with deterministic intensity function $\nu(t)$.

We denote by $\wb q$ the queue length process after insertion of the metaorder. Its dynamics are given by
\begin{equation}
\label{eq:bar_queue}
\wb q_t = q_0 + \wb L_t - \wb C_t - \wb N_t + L^o_t ,
\end{equation}
where $\wb L$, $\wb C$, and $\wb N$ are inhomogeneous Poisson processes with stochastic intensities
\begin{equation}
\wb \lambda_t^{L} = \lambda^L(\wb q_{t-}),
\qquad
\wb \lambda_t^{C} = \lambda^C(\wb q_{t-}), \qquad \wb \lambda_t^{N} = \lambda^N(\wb q_{t-}).
\end{equation}

We denote by
\[
\mathcal F_\infty = \sigma\bigl( (L_t, C_t, N_t, q_t),\, t \ge 0 \bigr)
\]
the filtration generated by the observed queue dynamics.
The objective of this section is to simulate trajectories of $(\bar{L}, \bar{C}, \bar{N}) | \mathcal{F}_\infty$. In other words, we want to simulate trajectories the shocked queue $\wb q$ conditionally on the observed trajectory of the original queue $q$.

We now describe the corresponding conditional simulation procedure.

\begin{proposition}
\label{prop:conditional_simulation}
The conditional simulation of $(\bar{L}, \bar{C}, \bar{N})$ given $\mathcal{F}_\infty$ is performed iteratively in time.

Assume that the processes have been simulated up to time $t$. 
Conditionally on the current state $\bar q_{t}$ and the filtration $\mathcal{F}_\infty$, the next jump time of the queue is determined as follows:

\begin{itemize}
    \item For $x \in \left \{ L, C , N\right \}$, we draw $U_j^{x} \sim \mathcal{U}([0,1])$ tied to each jump event $T_j^x$ of the process $x$.
    \item $\tau^L \sim \mathcal{E}((\lambda^L(\bar{q}_t) - \lambda^L(q_t))_+)$.
    \item $\tau^C \sim \mathcal{E}((\lambda^C(\bar{q}_t) - \lambda^C(q_t))_+)$.
    \item $\tau^N \sim \mathcal{E}((\lambda^N(\bar{q}_t) - \lambda^N(q_t))_+)$.
    \item $\tau^{L^o}= \inf \left \{s \geq 0: L^o_{s+t} = L^o_t+1\right \}$.
    \item $\bar{\tau}^L = \inf \left \{s \geq 0: \exists j \text{ such that } s = T_j^L - t, U_j^L \lambda^L(q_t) \leq \lambda^L(\bar{q}_t) \right \}$.
    \item $\bar{\tau}^C = \inf \left \{s \geq 0: \exists j \text{ such that }, s = T_j^C - t, U_j^C \lambda^C(q_t) \leq \lambda^C(\bar{q}_t) \right \}$.
    \item $\bar{\tau}^N = \inf \left \{s \geq 0: \exists j \text{ such that }, s = T_j^N - t, U_j^N \lambda^N(q_t) \leq \lambda^N(\bar{q}_t) \right \}$.
\end{itemize}

We then consider the next event time $\tau = \min \left \{ \tau^L, \tau^C, \tau^N, \tau^{L^o}, \bar{\tau}^L, \bar{\tau}^C, \bar{\tau}^N \right \}$ so that the next event occurs at time $t + \tau$. The nature of the event is determined as follows:
\begin{itemize}
    \item Limit order: it occurs if $\tau \in \left \{ \tau^L , \bar{\tau}^L, \tau^{L^o} \right \}$.
    \item Cancel order: it occurs if $\tau \in \left \{ \tau^C , \bar{\tau}^C \right \}$.
    \item Market order: it occurs if $\tau \in \left \{ \tau^N , \bar{\tau}^N \right \}$.
\end{itemize}
\end{proposition}
The proof is given in Appendix \ref{appendix:section_four_one}.

To illustrate our results, we fix a finite time horizon $T>0$ and consider the dynamics of all processes on the interval $[0,T]$. 
We assume that the metaorder arrival process $L^o$ is time-homogeneous, with constant intensity $\nu>0$. 
Figure \ref{fig:queue_with_limit} shows the conditional distribution of $\bar q$ under the specifications above.
Similarly, we consider an incoming metaorder $N^o$ represented by a time homogeneous Poisson process with intensity $\eta > 0$ executed exclusively through market orders. We perform the corresponding conditional simulation. The corresponding shocked queue dynamics are given by
\begin{equation*}
\wb q_t = q_0 + \wb L_t - \wb C_t - \wb N_t - N^o_t.
\end{equation*}
The results are shown in Figure~\ref{fig:queue_with_market}. 
 \begin{figure}[H]
        \centering
        \includegraphics[width=0.85\textwidth]{images/queue_values_plot_limit.pdf}
        \caption{Conditional simulation of $q$ knowing $\wb q$ in the presence of a limit metaorder. The intensities processes $L$, $C$, $N$, and $L^o$ are respectively $\lambda^L(x) = 100-0.275x$, $\lambda^C(x) = 2+0.125x$, $\lambda^N(x) = 0.105x$, and $\nu = 5$. The time horizon $T = 100$ seconds and the metaorder lasts for 80 seconds.}
        \label{fig:queue_with_limit}
    \end{figure}
     \begin{figure}[H]
        \centering
        \includegraphics[width=0.85\textwidth]{images/queue_values_plot_market.pdf}
        \caption{Conditional simulation of $q$ knowing $\wb q$ in the presence of an aggressive metaorder. The intensities processes $L$, $C$, $N$, and $L^o$ are respectively $\lambda^L(x) = 100-0.275x$, $\lambda^C(x) = 2+0.125x$, $\lambda^N(x) = 0.105x$, and $\eta = 5$. The time horizon $T = 100$ seconds and the metaorder lasts for 80 seconds.}
        \label{fig:queue_with_market}
    \end{figure}
\begin{remark}
Notice that the queues become indistinguishable shortly after the metaorder ends. 
We explain this in two steps:

\begin{itemize}
    \item Under the queue dynamics \eqref{eq:queue} and \eqref{eq:bar_queue}, it is shown in \cite{youssef2024passiveimpact} that 
    $\wb q_t \ge q_t$ (resp. $\wb q_t \le q_t$) for all $t \in [0,T]$ when executing $L^o$ (resp. $N^o$).
    
    \item Since $\lambda^L - \lambda^C - \lambda^N$ is decreasing, both queues are driven toward the same ergodic regime after the metaorder ends. 
    As they are governed by the same underlying point process, their difference vanishes, and they therefore coincide pathwise.
\end{itemize}
\end{remark}
\subsection{Market impact estimation}\label{subsec:impact_estimation}
Market impact captures the average price response to trading: buy orders induce upward price movements, while sell orders induce downward movements. 
Formally, it is defined as the difference between the price trajectory observed during the execution of a metaorder and the counterfactual trajectory that would have happened had the metaorder not been executed. 

Thus, by simulating conditional queue trajectories with and without a metaorder, our methodology provides direct access to the distribution of market impact. In particular, it enables two complementary conditional simulations:
\begin{itemize}
    \item \textbf{Market impact \textit{a priori}}: given an observed queue trajectory $q$, we simulate the conditional distribution of the shocked queue $\wb q$, allowing to test multiple trading strategies and quantify their impact.
    \item \textbf{Market impact \textit{a posteriori}}: given an observed queue trajectory $\wb q$ in the presence of a metaorder, we simulate the conditional distribution of the counterfactual queue process $q$, thereby quantifying the market impact generated by the realized metaorder strategy.
\end{itemize}

In the same spirit as \cite{Chahdi2024}, we  model the aggregated volumes available on the ask and bid side by $q^a$ and $q^b$ and we write
\begin{equation}
\label{eq:queue_dynamics}
q^a_t = L^{a}_t - C^{a}_t - N^{a}_t \quad \text{ and } \quad q^b_t = L^{b}_t - C^{b}_t - N^{b}_t
\end{equation}
where $L^a_t$ and $L^b_t$ are the number of limit orders on the ask and the bid side before time $t$, $C^a_t$ and $C^b_t$ are the number of cancel orders on the ask and the bid side before time $t$, and $N^a_t$ and $N^b_t$ are the number of buy and sell market orders before time $t$. We assume that $L^{a}$ and $L^{b}$ are two time-inhomogeneous Poisson processes whose intensities are given by 
\begin{equation*}
\lambda_t^{L,a} = \lambda^L(q_{t-}^{a})
\quad \text{ and } \quad 
\lambda_t^{L,b} = \lambda^L(q_{t-}^{b})
\end{equation*}
for some given decreasing function $\lambda^L$; and that $C^{a}$ and $C^{b}$ are two time-inhomogeneous Poisson processes whose intensities are given by 
\begin{equation*}
\lambda_t^{C,a} = \lambda^C(q_{t-}^{a})
\quad \text{ and } \quad 
\lambda_t^{C,b} = \lambda^C(q_{t-}^{b})
\end{equation*}
for some given increasing function $\lambda^C$.\\
Furthermore,  we assume that $N^a$ and $N^b$ are two independent Hawkes processes with the same baseline intensity $\mu \geq 0$ and self-exciting (non-negative) kernel $\varphi$ so that their intensities are given by
\begin{equation*}
\lambda^{a}_t = \mu + \int_0^{t-} \varphi(t-s) \, dN^{a}_s
\;\;\text{ and }\;\;
\lambda^{b}_t = \mu + \int_0^{t-} \varphi(t-s) \, dN^{b}_s.
\end{equation*}
In contrast with the model of Section \ref{subsec:lob_sim}, market orders in this setup are independent of the queue state at any time and depend solely on the history of market orders.

The authors in \cite{youssef2024passiveimpact} define the price as the limit of the anticipation of the upcoming market orders with a weight that depends on the limit order book state at the time of execution of the market order. This weight encodes the fact that it is harder for the price to move in one direction when the available volume in this side of the limit order book is higher. The price is then given by
\begin{equation}
\label{eq:price}
P_t = P_0 + \lim_{T\to\infty}
\E\bigg[\int_0^T \kappa(q^a_s) \, dN^a_s - \int_0^T \kappa(q^b_s) \, dN^b_s \,\bigg|\, \calF_t \bigg]
\end{equation}
for some function decreasing function $\kappa$. 
\subsubsection{Passive market impact}
We introduce a metaorder executed exclusively through limit orders on the ask side. The resulting queue process is similar to that given by \eqref{eq:bar_queue}
\begin{equation*}
\wb q^a_t = q_0^a + \wb L^{a}_t - \wb C^{a}_t - N^{a}_t + L^{o}_t
\end{equation*}
where $L^o$ denotes a non-homogeneous Poisson process representing the metaorder, $\wb L^{a}$ and $\wb C^{a}$ are Poisson inhomogeneous processes with stochastic intensities
\begin{equation*}
\wb \lambda_t^{L,a} = \lambda^L(\wb q_{t-}^{a})
\quad \text{ and } \quad 
\wb \lambda_t^{C,a} = \lambda^C(\wb q_{t-}^{a}).
\end{equation*}
The price process in presence of the metaorder is given by
\begin{equation}
\label{eq:price:lim:metaorder}
\wb P_t = P_0 + \lim_{T\to\infty}
\E\bigg[\int_0^T \kappa(\wb q^{a,t}_{s}) \, dN^a_s - \int_0^T \kappa(q^b_s) \, dN^b_s \,\bigg | \, \calF_t \bigg].
\end{equation}
where 
\begin{equation*}
\wb q^{a,t}_{s} = \wb L^{a,t}_{s} - \wb C^{a,t}_{s} - N^{a}_t + L_{s}^{o,t}
\quad \quad \text{ and } \quad \quad
 L_{s}^{o,t} = L^{o}_{s \wedge t}
\end{equation*}
The truncation at time $t$ reflects the fact that, while the market digests information from past order flow, it cannot anticipate the arrival of future metaorders. Consequently, when computing the anticipated future flow, the metaorder must be truncated at time $t$. 

The pathwise passive market impact is then defined as
\begin{equation}
\label{eq:mi}
MI_t = \wb P_t - P_t = \lim_{T\to\infty}
\E\bigg[\int_0^T \big( \kappa(\wb q^{a,t}_{s}) - \kappa(q^a_s) \big) \, dN^a_s \,\bigg | \, \calF_t \bigg].
\end{equation}

Our conditional simulation approach can be applied within this framework to quantify the impact of limit orders, or in other words, passive market impact. Although computing passive market impact is possible in this more general setting, it nevertheless requires a Monte Carlo simulation of the queue dynamics at each time step. We propose a much more efficient simulation method, under the assumption that the Hawkes processes have exponential kernels, and that $\kappa$ and $\lambda^L - \lambda^C$ are affine functions.
Similarly to \cite{youssef2024passiveimpact}, we consider
\begin{as}
    
\end{assumption}
\begin{equation*}
    \kappa(x) = d_\kappa + c_\kappa x, \qquad \text{and} \qquad \lambda^L(x) - \lambda^C(x) = c_\lambda x + d_\lambda,
\end{equation*}
where $c_\kappa < 0$, to guarantee the decreasing behavior of the impact function, and $c_\lambda < 0$ to reflect that $\lambda^L - \lambda^C$ is decreasing. The market impact in \eqref{eq:mi} becomes
\begin{equation*}
    MI_t = \lim_{T \to \infty} \E\bigg[\int_0^T \big( \kappa(\wb q^{a,t}_{s}) - \kappa(q^a_s) \big) \, dN^a_s \,\bigg | \, \calF_t \bigg] = c_{\kappa} \E\left[\int_{0}^{+\infty} (\bar{q}^{a,t}_s - q^a_s) dN^a_s \bigg| \mathcal{F}_t\right]
\end{equation*}
for $t > 0$.

The main approximation result is the following proposition.

\begin{proposition}
\label{prop:nonpredictablemartingale}
\begin{equation*}
    MI_t = c_{\kappa} \int_0^t (\bar{q}^{a,t}_s - q^{a}_s) dN^a_s + c_{\kappa} (\bar{q}^{a,t}_t - q^{a}_t)\int_t^{+\infty} e^{c_{\lambda}(s-t)} \E[\lambda^a_s|\mathcal{F}_t] ds.
\end{equation*}
\end{proposition}

We assume that $N^a$ admits a self-exciting kernel given by a sum of exponentials, so that its intensity is given by
 $\lambda^a_t = \mu + \int_0^t \varphi(t-s)dN^a_s$ where 
 \begin{equation}
 \label{eq:kernel_sum_exps}
 \varphi(t) = \sum_{i=1}^m \alpha_i e^{-\beta_i t}
 \end{equation}
 for $m \ge 1$ and some positive constants $(\alpha_i,\beta_i)_{i \in \{1,\cdots, m\}}$ satisfying $\norm{\phi}_1 := \sum_{i=0}^m \frac{\alpha_i}{\beta_i} < 1$.\\
 The Hawkes process becomes Markovian in the factors $R^i_t := \int_0^t \alpha_i e^{-\beta_i (t-s)} dN^a_s$, for $i \in \{1,\cdots,m\}$.\\
Under these dynamics, an explicit form for the market impact \eqref{eq:mi} is found and is given in the following theorem.
\begin{theorem}
    \label{thm:explicit_impact_shape}
    With coefficients $(x_i)_{i \in \{1,\cdots,m\}}$ and $C$ explicitely accessible, one has:
    \begin{equation*}
        MI_t = c_{\kappa} \int_0^t (\bar{q}^{a,t}_s - q^{a}_s) dN^a_s + c_{\kappa} (\bar{q}^{a,t}_t - q^{a}_t)(\sum_{i=1}^m x_i R^i_t + C)
    \end{equation*}
\end{theorem}

The following propositions are used to prove proposition \ref{prop:nonpredictablemartingale} and the theorem \ref{thm:explicit_impact_shape}. The proofs for this section are given in Appendix \ref{appendix:section_four_two}.

We have the following result on which proposition \ref{lem:queue_independence} relies.
\begin{lemma}
    \label{lem:queue_independence}
    For all $s > t$, we have
    \begin{equation*}
    \E[(\bar{q}^{a,t}_s - q^a_s) \lambda^a_s | \mathcal{F}_t] = \E[\bar{q}^{a,t}_s - q^a_s | \mathcal{F}_t] \E[\lambda^a_s | \mathcal{F}_t].
    \end{equation*}
Furthermore, under the affine consideration of $\lambda^L - \lambda^C$, we have 
\begin{equation*}
    \E[\bar{q}^{a,t}_s - q^a_s | \mathcal{F}_t] = (\bar{q}^a_t - q^a_t) e^{c_{\lambda}(s-t)}.
\end{equation*}
\end{lemma}

 To get now to the explicit formula, we solve explicitely the conditional expectation of the intensity process. \\
 The process $f_t(s) := \E_t[\lambda_s] = \E[\lambda^a_s|\mathcal{F}_t] $ solves the following Volterra equation:
    \begin{equation}
    f_t(s) = \mu + \sum_{i=1}^m R^i_t e^{-\beta_i (s-t)} + \int_t^s \left(\sum_{i=1}^m \alpha_i e^{-\beta_i(s-u)}\right) f_t(u) du
    \end{equation}
    The solution to this equation relies on finding the inverse of the convolution operator $\delta_0 - \sum_{i=1}^m \alpha_i e^{-\beta_i (\cdot)}$. We then have the following results:
\begin{lemma}
    \label{lem:propagator_formula}
    There exist positive constants $((\lambda_i, c_i))_{i=1}^m$ such that the inverse of the convolution operator $\delta_0 - \sum_{i=1}^m \alpha_i e^{-\beta_i (\cdot)}$ is $\delta_0 + \sum_{i=1}^m c_i e^{-\lambda_i (\cdot)}$.
\end{lemma}
\begin{lemma}
    \label{lem:closed_impact_term}
    Under the Hawkes kernel shape \eqref{eq:kernel_sum_exps}, we have
    \begin{align*}
    \int_{t}^{+\infty} e^{c_{\lambda}(s-t)} \E_t[\lambda^a_s]ds = &\sum_{i}^{m} R^i_t \left( \frac{1}{ \beta_i - c_{\lambda}} + \sum_{j=1}^m \frac{c_j}{\lambda_j - \beta_i}\left(\frac{1}{\beta_i - c_\lambda} - \frac{1}{\lambda_j - c_\lambda}\right)\right) \\
    &+ \mu \left(1 -  \frac{1}{c_\lambda} + \sum_{j=1}^m \frac{c_j}{\lambda_j} \left(\frac{1}{c_\lambda - \lambda_j} - \frac{1}{c_\lambda}\right)\right).
    \end{align*}
    
    In particular, the representation is affine in the Markovian factors and we identity the coefficients $C$ and $(x_i)_{i \in \{ 1, \cdots, N\}}$ of the main statement in the formula above.
\end{lemma}
\begin{remark}
    Efficient simulation of market impact trajectories becomes possible through closed-form expressions, with no dependence on anticipated future order flow. From a technical standpoint, power-law kernels are better suited to replicate the long-memory behavior of trading activity observed in financial markets. In this case, various approximation techniques allow power-law kernels to be approximated by sums of exponentials, see \cite[]{}, thereby making market impact simulation both easier and scalable within this modeling framework.
\end{remark} 
In particular, Figure \ref{} 

\subsubsection{Aggressive market impact}
In the following, we consider a simplified version of the model in \eqref{eq:price}, where the price is given by
\begin{equation}
\label{eq:simplified}
  P_t = P_0 + \int_0^t \kappa(q^a_s) \xi(t-s) dN^a_s - \int_0^t \kappa(q^b_s) \xi(t-s) dN^b_s.
\end{equation}
Equation~\eqref{eq:simplified} can be interpreted as a reduced version of the model \eqref{eq:price}, in which two key mechanisms are deliberately omitted. 
First, market orders are the sole contributors to the quadratic variation of the price in \eqref{eq:simplified}. 
The direct effects of limit and cancellation orders are removed and only influence the price indirectly through the cumulative queue volumes $q^a$ and $q^b$. 
Second, the conditional expectation appearing in the price definition \eqref{eq:price} is suppressed and replaced by a closed-form expression analogous to the constant-$\kappa$ specification studied in \cite{jaisson2015market}.

Under this specification, we introduce a buy metaorder, denoted $N^o$, executed through market orders on the ask side. The market impact in this case is given by
\begin{equation}
    MI_t = \int_0^t \left(\kappa(\wb q^{a}_s) - \kappa(q^a_s)\right) \xi(t-s) dN^a_s + \int_0^t \kappa(\wb q^{a}_s) \xi(t-s) dN^o_s.
\end{equation}
where
\begin{equation*}
\wb q^a_t = q_0^a + \wb L^{a}_t - \wb C^{a}_t - N^{a}_t - N^{o}_t.
\end{equation*}
The simulation methodology is identical to that of the previous section. We first simulate an initial queue trajectory $q^a$ (observed realization) in the absence of a metaorder, together with a single realization of the metaorder process $N^o$, which is kept fixed throughout the experiment. Conditional on the observed trajectory $q^a$, we then generate multiple realizations of the conditional queue process $\wb q^a$, all incorporating the execution of the same metaorder, and compute the corresponding market impact. Figure~\ref{} illustrates the resulting distribution of market impact.

\bibliographystyle{plainnat}
\bibliography{references}

\appendix
\section{Proofs of Section \ref{subsec:main_results}} \label{appendix:section_three_tree}

\closuremeasurability*

\begin{proof}[Proof of Lemma \ref{lem:closuremeasurability}]
    \leavevmode \\
    We will start by showing that the function $G: \mathbb{L} \to C(E)$ given by $g \to cl(\{ (s,k,z) \in E | z \leq g_k(s)\})$ is measurable.\\
    Let $O = (a,b) \times \{k\} \times (c,d)$ be an elementary open of $E$ and $C_O$ the corresponding open of $C(E)$.\\
    We have the following equalities:
    \begin{align*}
        G^{-1}(C_O) &= \{ g \in \mathbb{L}| B(g) \cap O \neq \varnothing \} \\
        &= \{ g \in \mathbb{L}|A(g) \cap O \neq \varnothing \} \text{ as } B(g) = cl(A(g)). \\
        &= \{ g \in \mathbb{L}| \exists (t,z) \in (a,b) \times (c,d), z \leq g_k(t) \} \\
        &= \{ g \in \mathbb{L} | \sup_{(a,b)} g_k > c\} \\
        &= \{ g \in \mathbb{L} | \sup_{(a,b) \cap \mathbb{Q}} g_k > c \} \text{ as } \sup_{(a,b)} g_k = \sup_{(a,b) \cap \mathbb{Q}} g_k \text{ for } g \in \mathbb{L}.\\
        G^{-1}(C_O)&= \cup_{t \in (a,b) \cap \mathbb{Q}} \{ g \in \mathbb{L} | g_k(t) > c \}
    \end{align*}
    This concludes that $G^{-1}(C_O)$ is measurable as $\{ g \in \mathbb{L} | g_k(t) > c \}$ is measurable for the Kolmogorov $\sigma$-algebra and this is a countable union of measurable sets.\\
    Now, $B(\pi)$ is a well-defined random set by composition of measurable maps as follows:
    \[
    \begin{tikzcd}
    B(\pi) \colon (\Omega, \mathbb{P}) \arrow[r, "\pi"] & \mathbb{M}(E) \arrow[r, "\lambda"] & \mathbb{L} \arrow[r, "G"] & C(E)
    \end{tikzcd}
    \]\\
    To conclude on the compactness, our assumption that  for all $\mu \in \mathbb{M}(E)$, $\sup_{t \in [0,T]} \sup_{i=1}^d \lambda^{\mu, i}_t < \infty$ imples that $A(\mu) \subset [0,T] \times \{1,\cdots,d\} \times [0, \sup_{t \in [0,T]} \sup_{i=1}^d \lambda^{\mu, i}_t]$ and the latter is compact. As such, the closure $B(\mu)$ is compact as a closed subset of a compact.

    This concludes that $B(\pi)$ is a compact random set.
\end{proof}

\stability*

We recall the following inductive property of the processes $\lambda^\mu$ and $N^\mu$ that will be used in the proof.
\begin{align*}
    \lambda^\mu_t &= f(j_{t,*}(N^\mu), t) \\
    N^\mu_t &= \int_E e_k \ind_{\{z \le \lambda^{\mu,k}_s\}} \ind_{ \{s \leq t \} } \, \mu(\diff s, \diff k ,\diff z)
\end{align*}

\begin{proof}[Proof of Lemma \ref{lem:stability}]
    \leavevmode \\
    For conciseness of the proof, we extend $\lambda$, $f$ and $N$ to $]-\infty, 0 [$ by $\lambda^\mu_s = \lambda^\mu_0$, $f(-,s) = 0$ and $N^\mu_s = 0$ for $s<0$. This extension to $]-\infty,0[$ does not change the model on $[0,T]$. We will denote in the following proof $\tilde{\mu} := \mu_{|B(\mu)} \cup \mu'_{|B(\mu)^c}$.
    
    As $N^\mu_t = \int_{A(\mu) \cap ([0,t] \times \{ 1,\cdots, d \} \times \R_+)} e_k \mu(ds, dk, dz)$ for $\mu$ and $\tilde{\mu}$ combined with the compactness of $A(\mu)$ and $A(\tilde{\mu})$ implies that the processes $(\lambda^\mu, N^\mu)$ and $(\lambda^{\tilde{\mu}}, N^{\tilde{\mu}})$ are determined by the restriction of $\mu \cup \mu'$ to a large enough compact K.\\
    Let $S$ be the union of $\left \{ -1 \right \}$ and the projection of the support of $(\mu \cup \mu')_{|K}$ to the time axis. \\
    The latter set is now a subset of $\R_+$ without accumulation as its the continuous projection from a compact without accumulation points. \\
    We can now write $S$ as a strictly increasing sequence of times and use the induction principle on it. \\
    We will show by induction that for all $t \in S$, $$\mathcal{P}(t): A(\tilde{\mu})_{|[0,t]} = A(\mu)_{|[0,t]} \text{ and } N^{\tilde{\mu}}_{|]-\infty,t]} = N^{\mu}_{|]-\infty,t]}.$$

    \begin{itemize}
        \item \underline{Initialization:} $\mathcal{P}(-1)$ holds.
        
        As -1 is not in the support of $\mu$ and $\tilde{\mu}$, $\mathcal{P}(-1)$ holds as:
        $$N^\mu_{-1} = N^{\tilde{\mu}}_{-1} = 0 \text{ and as } \lambda_{-1}^\mu = f(0, -1) = \lambda_{-1}^{\tilde{\mu}}$$

        \item \underline{Induction:}  Assume now that $\mathcal{P}(t_0)$ holds, and we want to show $\mathcal{P}(t)$ holds for the next event $t$ of $\pi \cup \pi'$ strictly larger than $t_0$.
        
        \begin{itemize}
            \item \underline{$N^\mu = N^\tilde{\mu}$ on $]t_0, t[$:}
            
            The following equalities hold as $\mu$ and $\mu'$ have no jumps in that interval inside K:
            \begin{align*}
            N_s^\mu &= N_{t_0}^\mu + \int_{t_0+}^s \int_{\{1,\cdots, d\} \times \R_+} e_k \ind_{\lp z \leq \lambda_u^\mu\rp} \mu(\diff u, \diff k, \diff z) \\
            &= N_{t_0}^{\tilde{\mu}} + \int_{t_0+}^s \int_{\{ 1,\cdots, d \} \times \R_+} e_k \ind_{\lp z \leq \lambda_u^{\tilde{\mu}} \rp} \tilde{\mu} (du, dk, dz) \\ 
            N_s^\mu &= N_s^{\tilde{\mu}}
            \end{align*}
            
            \item \underline{$A(\mu)_{|[0,t]} = A(\tilde{\mu})_{|[0,t]}$}:
            
            This now implies that $N^\mu_{|]-\infty, t[} = N^{\tilde{\mu}}_{|]-\infty,t[}$, which is equivalent to $j_{s,*}(N^\mu) = j_{s,*}(N^{\tilde{\mu}})$ for all $s\leq t$ and can conclude that: $$\forall s \in ]-\infty,t], \lambda^\mu_s = f(j_{s,*}(N^\mu), s) = f(j_{s,*}(N^{\tilde{\mu}}), s) = \lambda^{\tilde{\mu}}_s.$$
            
            \item \underline{Equality of $N^\mu$ and $N^\tilde{\mu}$ at t:}
            
            As $\lambda^\mu$ is a càglàd function, one has the following description of the closure of the hypograph: $$\operatorname{cl}(\left \{ (s,k,z) \in [0,T] \times \{ 1,\cdots, d \} \times \R_+ | z \leq \lambda^{\mu,k}_s \right \}) = \left \{ (s,k,z) \in [0,T] \times \{ 1,\cdots, d \} \times \R_+ | z \leq \max(\lambda^{\mu,k}_s,\lambda^{\mu,k}_{s+}) \right \}.$$
            
            This identity comes from the fact that in dimension one, for a càglàd function $g$ and a sequence $(z_n, s_n)$ of points of the hypograph converging to a pair $(z,s)$, the inequality $z_n \leq g(s_n)$ translates to the $\limsup$ as $z \leq \limsup_n f(s_n) \leq \max(f(s), f(s+))$ and this bound can conversely be reached explicitely.
            
            Let us now write $N^{\tilde{\mu}}_{t} = N^{\tilde{\mu}}_{t-} + \sum_{k=1}^d e_k \tilde{\mu}(\left \{ t \right \} \times \left \{ k \right \} \times [0, \lambda^{\tilde{\mu},k}_t])$ and $\lambda^{\tilde{\mu}}_t = \lambda^\mu_t$.
            
            The second term splits in two terms: 
            
            $$\tilde{\mu}(\left \{ t \right \} \times \left \{ k \right \} \times [0, \lambda^{\tilde{\mu},k}_t]) = \mu((\left \{ t \right \} \times \left \{ k \right \} \times [0, \lambda^{\tilde{\mu}}_t]) \cap B(\mu)) + \mu'((\left \{ t \right \} \times \left \{ k \right \} \times [0, \lambda^{\tilde{\mu}}_t]) \cap B(\mu)^c).$$
            
            They can be controled explicitely thanks to the precise description of the closure of the hypograph:
            \begin{align*}
                (\left \{ t \right \} \times \left \{ k \right \} \times [0, \lambda^{\tilde{\mu}}_t]) \cap B(\mu) &= \left \{ t \right \} \times \left \{ k \right \} \times ([0, \lambda^{\tilde{\mu}}_t] \cap [0, \max(\lambda^\mu_t, \lambda^\mu_{t+})]) \\
                (\left \{ t \right \} \times \left \{ k \right \} \times [0, \lambda^{\tilde{\mu}}_t]) \cap B(\mu)^c &= \left \{ t \right \} \times \left \{ k \right \} \times ([0, \lambda^{\tilde{\mu}}_t] \cap ]\max(\lambda^\mu_t, \lambda^\mu_{t+}), +\infty[)
            \end{align*}
            
            This simplifies then to the following equations:
            \begin{align*}
                (\left \{ t \right \} \times \left \{ k \right \} \times [0, \lambda^{\tilde{\mu}}_t]) \cap B(\mu) &= \left \{ t \right \} \times \left \{ k \right \} \times [0, \lambda^\mu_t] \\
                (\left \{ t \right \} \times \left \{ k \right \} \times [0, \lambda^{\tilde{\mu}}_t]) \cap B(\mu)^c &= \emptyset
            \end{align*}
            
            We can then conclude that $\tilde{\mu}(\left \{ t \right \} \times \left \{ k \right \} \times [0, \lambda^{\tilde{\mu}, k}_t]) = \mu(\left \{ t \right \} \times \left \{ k \right \} \times [0, \lambda^{\mu,k}_t])$, which yields $N^{\tilde{\mu}}_t = N^\mu_t$.
            
        \end{itemize}
    \item By induction, we can conclude that $\mathcal{P}(t)$ holds for all $t \in S$.
    
    Denoting $t = \sup_{s \in S} s$, the first two points of the induction step allow to conclude now that the equality holds on $]t, T]$ as the measures no longer have support intersecting K and as such, one concludes that $\mathcal{P}(T)$ holds.
    \end{itemize}
\end{proof}

\stoppingset*

\begin{proof}[Proof of Lemma \ref{lem:stoppingset}]
    \leavevmode \\
    Let $L$ be a compact in $E$, we will first show that $\{ m \in \mathbb{M}(E) | B(m) \subset L \} = \{ m \in \mathbb{M}(E) | B(m_{|L}) \subset L \}$.\\
    Let $m \in \mathbb{M}(E)$ with $B(m) \subset L$ and with $(\mu,\mu') := (m,m_{|L})$, Lemma \ref{lem:stability} yields:
    \begin{align*}
        B(m) &= B(m_{|B(m)} \cup m_{|L \cap B(m)^c}) \\
        &= B(m_{B(m) \cup (L \cap B(m)^c)}) \\
        &= B(m_{|L})
    \end{align*}
    In particular, $B(m_{|L}) \subset L$ and this concludes on the first inclusion.\\
    Conversely, let $m \in \mathbb{M}(E)$ be such that $B(m_{|L}) \subset L$.\\
    Lemma \ref{lem:stability} yields $B(m_{|L}) = B((m_{|L})_{|B(m_{|L})} \cup m_{|B(m_{|L})^c})$ with $(\mu, \mu') := (m_{|L}, m)$.

    Since $B(\mu_{|L}) \subset L$, $(\mu_{|L})_{|B(\mu_{|L})} = \mu_{|B(\mu_{|L})}$ and we get:
    \begin{align*}
        B(\mu_{|L}) &= B((\mu_{|L})_{|B(\mu_{|L})} \cup \mu_{|B(\mu_{|L})^c}) \\
        &= B(\mu_{|B(\mu_{|L})} \cup  \mu_{|B(\mu_{|L})^c}) \\
        &= B(\mu)
    \end{align*}
    This concludes the equality and as $(B(\pi_{|L}) \subset L)$ is $\sigma(\pi_{|L})$-measurable, this concludes that $(B(\pi) \subset L)$ is $\sigma(\pi_{|L})$-measurable.

    This holds for every compact $L$ of $E$ and $B(\pi)$ is thus a stopping set with respect to $\pi$.
\end{proof}

We now finish by applying the strong Markov property to the stopping set $B(\pi)$.

\conditionallaw*

\begin{proof}[Proof of Theorem \ref{thm:conditionallaw}]
        \leavevmode \\
        The first part will be to show the conditional independence and will be split in the following steps:
        \begin{itemize}
            \item For $\pi, \pi'$ i.i.d., we show $\mathcal{L}((\pi_{|B(\pi)}, \pi_{|B(\pi)^c})) = \mathcal{L}((\pi_{|B(\pi)}, \pi'_{|B(\pi)^c}))$:\\
            Let $f,g: \mathbb{M}(E) \to \R_+$ measurable,\\
            Theorem \ref{thm:strongmarkov} applied to $\phi := (\mu \in \mathbb{M}(E) \to f(\mu_{|B(\mu)}) g(\mu_{|B(\mu)^c}) \in \R_+)$ yields:
            \begin{align*}
                \E[\phi(\pi)] &= \E[\phi(\pi_{|B(\pi)} \cup \pi'_{|B(\pi)^c})] \\
                \E[f(\pi_{|B(\pi)})g(\pi_{|B(\pi)^c})] &= \E[f((\pi_{|B(\pi)} \cup \pi'_{|B(\pi)^c})_{|B(\pi_{|B(\pi)} \cup \pi'_{|B(\pi)^c})}) g((\pi_{|B(\pi)} \cup \pi'_{|B(\pi)^c})_{|B(\pi_{|B(\pi)} \cup \pi'_{|B(\pi)^c})^c})] \\
                \E[f(\pi_{|B(\pi)})g(\pi_{|B(\pi)^c})] &= \E[f((\pi_{|B(\pi)} \cup \pi'_{|B(\pi)^c})_{|B(\pi)}) g((\pi_{|B(\pi)} \cup \pi'_{|B(\pi)^c})_{|B(\pi)^c})] \text{ as } B(\pi) = B(\pi_{|B(\pi)} \cup \pi'_{|B(\pi)^c}). \\
                \E[f(\pi_{|B(\pi)})g(\pi_{|B(\pi)^c})] &= \E[f(\pi_{|B(\pi)}) g(\pi'_{|B(\pi)^c})]
            \end{align*}

            \item $\pi_{|B(\pi)}$ and $\pi_{|B(\pi)^c}$ are independent conditionally on $B(\pi)$:\\
            We will do this by showing that for $f,g: \mathbb{M}(E) \to \R_+$: $$\forall Z \, \sigma(B(\pi))\text{-measurable}, \, \E[Zf(\pi_{|B(\pi)})g(\pi_{|B(\pi)^c})] = \E[Zf(\pi_{|B(\pi)}) g(\pi'_{|B(\pi)^c})].$$

            $\sigma(B(\pi))$ is generated by the events $\{ (B(\pi) \in C_O) | O \text{ open of } E$, we are then reduced to showing the equality for $Z = 1_{B(\pi) \in C_O}$.

            As $h := (\mu \in \mathbb{M}(E) \to 1_{B(\mu) \in C_O} \in \R_+)$ is measurable, we apply the first point to $hf$ and $g$:
            \begin{align*}
                \E[h(\pi_{|B(\pi)})f(\pi_{|B(\pi)})g(\pi_{|B(\pi)^c})] &= \E[h(\pi_{|B(\pi)})f(\pi_{|B(\pi)}) g(\pi'_{|B(\pi)^c})] \\
                \E[1_{B(\pi_{|B(\pi)}) \in C_O}f(\pi_{|B(\pi)})g(\pi_{|B(\pi)^c})] &= \E[1_{B(\pi_{|B(\pi)}) \in C_O}f(\pi_{|B(\pi)})g(\pi_{|B(\pi)^c})] \\
                \E[1_{B(\pi \in C_O)}f(\pi_{|B(\pi)})g(\pi_{|B(\pi)^c})] &= \E[1_{B(\pi \in C_O)}f(\pi_{|B(\pi)})g(\pi_{|B(\pi)^c})] \text{ as } B(\pi_{|B(\pi)}) = B(\pi).
            \end{align*}

            This provides the conditional expectation $\E[f(\pi_{|B(\pi)})g(\pi_{|B(\pi)^c})|\sigma(B(\pi))] = \E[f(\pi_{|B(\pi)}) g(\pi'_{|B(\pi)^c})|\sigma(B(\pi))]$.

            $\pi_{|B(\pi)}$ and $\pi_{|B(\pi)^c}$ are thus independent conditionally on $B(\pi)$ and as $\pi, \pi'$ are independent, $\mathcal{L}(\pi_{|B(\pi)^c}|B(\pi)) = \mathcal{L}(\pi'_{|B(\pi)^c}|B(\pi)) = \mathcal{P}oisson(1_{B(\pi)^c} \nu)$.

            \item Refining the statement conditionally on $\mathcal{F}^{N^\pi}_\infty$:\\
            We have $\sigma(B(\pi)) \subset \mathcal{F}^{N^\pi}_\infty \subset \sigma(\pi_{|B(\pi)})$ as $\lambda^\pi$ is a functional of $N^\pi$ and $N^\pi$ is a functional of $\pi_{|B(\pi)}$.

            The previous part tells us that $\pi_{|B(\pi)^c} \perp \sigma(\pi_{|B(\pi)})$ conditionally on $\sigma(B(\pi))$ and as $\mathcal{F}^{N^\pi}_\infty \subset \sigma(\pi_{|B(\pi)})$, we deduce $\pi_{|B(\pi)^c} \perp \mathcal{F}^{N^\pi}_\infty$ conditionally on $\sigma(B(\pi))$.

            As $\sigma(B(\pi))\subset \mathcal{F}^{N^\pi}_\infty$, this gives $\mathcal{L}(\pi_{|B(\pi)^c}|\mathcal{F}^{N^\pi}_\infty) = \mathcal{L}(\pi_{|B(\pi)^c}|B(\pi)) = \mathcal{P}oisson(1_{B(\pi)^c} \nu)$.

            We can then conclude on the conditional independence:\\
            Let $f,g: \mathbb{M}(E) \to \R_+$ measurable,
            \begin{align*}
                \E[f(\pi_{|B(\pi)}) g(\pi_{|B(\pi)^c}) | \mathcal{F}^{N^\pi}_\infty] &= \E[f(\pi_{|B(\pi)}) \E[g(\pi_{|B(\pi)^c})|\sigma(\pi_{|B(\pi)})]|\mathcal{F}^{N^\pi}_\infty] \\
                &= \E[f(\pi_{|B(\pi)}) \E[g(\pi_{|B(\pi)^c})|B(\pi)]|\mathcal{F}^{N^\pi}_\infty] \\
                &= \E[f(\pi_{|B(\pi)})|\mathcal{F}^{N^\pi}_\infty] \E[\E[g(\pi_{|B(\pi)^c})|B(\pi)]|\mathcal{F}^{N^\pi}_\infty] \\
                \E[f(\pi_{|B(\pi)}) g(\pi_{|B(\pi)^c}) | \mathcal{F}^{N^\pi}_\infty] &= \E[f(\pi_{|B(\pi)})|\mathcal{F}^{N^\pi}_\infty] \E[g(\pi_{|B(\pi)^c})|\mathcal{F}^{N^\pi}_\infty]
            \end{align*}
        \end{itemize}

        Finally, we describe the conditional law of the marks of $\pi_{|A(\pi)}$, let $\pi_{|A(\pi)} = \sum_{n \in I} \delta_{(T_n, k_n, Z_n)}$.

        For each $n \in I$, we have seen a jump of $N$ at $T_n$ in the dimension $k_n$.
    
        As such, this corresponds to the following condition on $Z_n$: 
        $$\mathcal{L}(Z_n| \mathcal{F}^{N^\pi}_\infty) = \mathcal{L}(Z_n|Z_n \leq \lambda^{\pi, k_n}_{T_n}) = \mathcal{U}([0, \lambda^{\pi, k_n}_{T_n}])$$
        as $\lambda^{\pi, k_n}_{T_n}$ is $\sigma(T_n) \cup \sigma((N_u)_{u<T_n})$-measurable, $Z_n \perp \sigma(T_n) \cup \sigma((N_u)_{u<T_n})$ and $Z_n$ drawn from the Lebesgue measure on $\R_+$.
\end{proof}
\section{Proofs of Section \ref{subsec:impact_estimation}}
\label{appendix:section_four_two}

\begin{proof}[Proof of Lemma \ref{lem:queue_independence}]
    \leavevmode \\
    According to the proof of Lemma B.4 in \cite{youssef2024passiveimpact}, there exists independent Poisson random measures $\tilde{\pi}$ and $\pi^N$ such that $(\bar{q}^{a,t}_s - q^a_s) = (\bar{q}^{a,t}_t - q^a_t) - \int_{t+}^s \int_{\R_+} 1_{z \leq -c_\lambda (\bar{q}^{a,t}_{u-} - q^a_{u-})} \tilde{\pi}(du, dz)$ and $N^a_s = \int_0^s \int_{\R_+} 1_{z \leq \lambda^a_u} \pi^N(du, dz)$ with $\lambda^a_s = \mu + \int_0^{s-} \phi(s-u) dN^a_u$.

    Denoting $U_s := \bar{q}^{a,t}_s - q^a_s$, can then write that:
    $$\mathcal{L}((U_{s}, \lambda^a_{s}) |\mathcal{F}_{t}) = \mathcal{L}((U_{s}, \lambda^a_{s}) |\sigma(U_{t}) \cup \sigma((N^a_{u})_{u \leq t}))$$
    \begin{align}
    	\mathbb{E}[U_{s} \lambda^a_{s} |\sigma(U_{t}) \cup \sigma((N^a_{u})_{u \leq t})] &= \mathbb{E}[\mathbb{E}[U_{s} \lambda^a_{s}|\sigma(U_{t}) \cup \sigma((N^a_{u})_{u \leq t}) \cup \sigma((N^a_{u})_{t < u \leq s})]|\sigma(U_{t}) \cup \sigma((N^a_{u})_{u \leq t})] \nonumber \\
    	&= \mathbb{E}[\lambda^a_{s} \mathbb{E}[U_{s}|\sigma(U_{t}) \cup \sigma((N^a_{u})_{u \leq s})] |\sigma(U_{t}) \cup \sigma((N^a_{u})_{u \leq t})] \nonumber
    \end{align}
    By Prop 11.10 of Le Gall, $\mathbb{E}[U_{s}|\sigma(U_{t}) \cup \sigma((N^a_{u})_{u \leq s})] = \mathbb{E}[U_{s}|\sigma(U_{t})] \text{ as } \sigma(U_{s})\cup \sigma(U_{t}) \perp \sigma((N^a_{u})_{u \leq s})$. 
    Can then conclude $\mathbb{E}[U_{s} \lambda_{s} |\mathcal{F}_{t}] = \mathbb{E}[\lambda_{s} \mathbb{E}[U_{s}|\mathcal{F}_{t}]|\mathcal{F}_{t}] = \mathbb{E}[U_{s} |\mathcal{F}_{t}] \mathbb{E}[\lambda_{s} |\mathcal{F}_{t}]$.

    Moreover, as $U_s = U_t - \int_{t+}^s \int_{\R_+} 1_{z \leq -c_\lambda U_{u-}} \tilde{\pi}(du, dz)$, the expectation solves $\E[U_s|\mathcal{F}_t] = U_t + c_\lambda \int_{t}^s \E[U_u|\mathcal{F}_t] du$ with unique solution $\E[U_s|\mathcal{F}_t] = U_t e^{c_\lambda (s-t)}$.
\end{proof}

In order to apply lemma \ref{lem:queue_independence} to lemma \ref{prop:nonpredictablemartingale}, we need to circumvent the fact that the Stieltjes integral is not made with respect to a predictable process.

\begin{proof}[Proof of Proposition \ref{prop:nonpredictablemartingale}]
    \leavevmode \\
    As $\bar{q}^{a,t}_s - q^a_s = (\bar{L}^{a,t}_s - L^a_s) + (\bar{C}^{a,t}_s - C^a_s) + L^o_s$, this is a point process driven by Poisson measures independent from the one driving $N^a$, this yields:
    $$\sum_{s\geq 0} \Delta(\bar{q}^{a,t} - q^a)_s \Delta N_s = 0 \text{ and } \int_0^s (\bar{q}^{a,t}_u - q^a_u)dN^a_u = \int_0^s (\bar{q}^{a,t}_{u-} - q^a_{u-})dN^a_u \text{ almost surely.}$$

    One then has $\int_{t+}^{+\infty} (\bar{q}^{a,t}_s - q^a_s) dN^a_s = \int_{t+}^{+\infty} (\bar{q}^{a,t}_{s-} - q^a_{s-}) dN^a_s$, taking the expectation yields: 
    $$\E\bigg[ \int_{t+}^{+\infty} (\bar{q}^{a,t}_s - q^a_s) dN^a_s  \bigg| \mathcal{F}_t\bigg] = \int_t^\infty \mathbb{E}[(\bar{q}^{a,t}_{s-} - q^a_{s-})\lambda^a_s]ds.$$

    As $\{ s \in \R_+ | \bar{q}^{a,t}_s - q^a_s \neq \bar{q}^{a,t}_{s-} - q^a_{s-} \}$ has zero Lebesgue measure, 
    $$\int_t^\infty \mathbb{E}[(\bar{q}^{a,t}_{s-} - q^a_{s-})\lambda^a_s]ds = \int_t^\infty \mathbb{E}[(\bar{q}^{a,t}_{s} - q^a_{s})\lambda^a_s]ds \text{ almost surely.}$$

    Thanks to lemma \ref{lem:queue_independence}, we have $\E[(\bar{q}^{a,t}_s - q^a_s)\lambda^a_s |\mathcal{F}_t] = (\bar{q}^a_t - q^a_t) e^{c_\lambda (s-t)}$ for $s \geq t$ and can conclude that:
    \begin{equation*}
    MI_t = c_{\kappa} \int_0^t (\bar{q}^{a,t}_s - q^{a}_s) dN^a_s + c_{\kappa} (\bar{q}^{a,t}_t - q^{a}_t)\int_t^{+\infty} e^{c_{\lambda}(s-t)} \E[\lambda^a_s|\mathcal{F}_t] ds.
    \end{equation*}
\end{proof}

\begin{proof}[Proof of Lemma \ref{lem:propagator_formula}]
    \leavevmode \\
    Working with Laplace transforms on $\R_+^*$, $F := \mathcal{L}\Bigl( (\delta_0 - \sum_{i=1}^m \alpha_i e^{-\beta_i (-)})^{-1} \Bigr) = \frac{1}{1 - \sum_{i=1}^m \frac{\alpha_i}{x + \beta_i}}$.
    
    We can rewrite $F$ as $F = 1 + \frac{P}{Q-P}$ with $deg(Q) > deg(P)$ and thus $deg(P) < deg(Q-P)$. 
    If we assume that $Q-P$ admits exactly $m$ distinct real roots $(\lambda_i)_{i=1}^m$, there exists then a decomposition $F = 1 + \sum_{i=1}^m \frac{c_i}{x + \lambda_i}$.
    
    We will now prove that $Q-P$ admits exactly $N$ distinct strictly negative roots.
    
    The roots are exactly the roots of the function $R(x) = 1 - \sum_{i=1}^{m} \frac{\alpha_i}{x+\beta_i}$ and we assume that the $\beta_i$ are ordered here.
    
    For each $i \in \{1,\cdots,m-1\}$, $\lim_{x \to -\beta_i^-}R(x) = +\infty$ and $\lim_{x \to -\beta_{i+1}^+} R(x) = -\infty$. By intermediate value theorem, we can conclude that there exists one zero $\lambda_{i+1} \in ]-\beta_{i+1}, -\beta_i[$ of $R$.
    
    Moreover, as $R$ is strictly increasing on $]-\beta_1, +\infty[$ with $\lim_{x \to -\beta_1^+}R(x) = -\infty$ and $R(0) = 1 - \sum_{i=1}^m \frac{\alpha_i}{\beta_i} > 0$, the intermediate value theorem provides the $m-th$ root $\lambda_1 \in ]-\beta_1, 0[$ of $R$.
    
    It remains to prove that the corresponding coefficients $c_i$ are all positive to conclude.
    
    In this decomposition here, one has $c_j = \frac{1}{R'(\lambda_j)}$ and as $R' = \sum_{i=1}^m \frac{\alpha_i}{(x+\beta_i)^2}$, this concludes on the positivity.
\end{proof}

\begin{proof}[Proof of Lemma \ref{lem:closed_impact_term}]
    With the notation $f_t(s) = 0$ for $s<t$ and $e^{-\beta (-)}$ be zero on negative times, one now has:
    $$(\delta_0 - \sum_{i=1}^m \alpha_i e^{-\beta_i (-)}) \ast f_t(s) = \sum_{i=1}^m R^i_t e^{-\beta_i(s-t)} + \mu$$

    
    Using the previous proposition, it implies:
    \begin{align*}
        f_t(s) &= \sum_{i=1}^m R^i_t \Bigl( e^{-\beta_i(s-t)} + \sum_{j=1}^m c_j \int_t^s e^{-\lambda_j (s-u)} e^{-\beta_i (u-t)} du \Bigr) + \mu \Bigl( 1  + \sum_{j=1}^m \frac{c_j}{\lambda_j}(1 - e^{-\lambda_j (s-t)})\Bigr).
    \end{align*}

    and moreover:
    \begin{align*}
    \int_t^s e^{-\lambda_j(s-u)} e^{-\beta_i(u-t)} du &= e^{\beta_i t - \lambda_j s} \int_t^s e^{u(\lambda_j - \beta_i)} du \\
    &= e^{\beta_i t - \lambda_j s} \frac{e^{s(\lambda_j - \beta_i)} - e^{t(\lambda_j - \beta_i)}}{\lambda_j - \beta_i} \\
    &= \frac{e^{-\beta_i (s-t)} - e^{-\lambda_j(s-t)}}{\lambda_j - \beta_i}
    \end{align*}
    
    Then, we can compute the integral of interest to us:
    \begin{align*}
        \int_t^{+\infty} e^{-c_{\lambda}(s-t)} \mathbb{E}_t[\lambda^a_s] ds 
        &= \begin{multlined}[t]
            \sum_{i=1}^m R^i_t \int_0^{+\infty} \Bigl(e^{-(c_\lambda + \beta_i) u} + \sum_{j=1}^m c_j \frac{e^{-\beta_i u} - e^{-\lambda_j u}}{\lambda_j - \beta_i} e^{-c_\lambda u}\Bigr) du \\
            + \mu \int_{t}^{\infty} e^{-c_\lambda (s-t)} \bigl(1 + \sum_{j=1}^N \frac{c_j}{\lambda_j}(1 - e^{-\lambda_j (s-t)})\bigr) ds
           \end{multlined} \\
        &= \sum_{i=1}^m R^i_t \Bigl(\frac{1}{c_\lambda + \beta_i} + \sum_{j=1}^m \frac{c_j}{\lambda_j - \beta_i} (\frac{1}{c_\lambda + \beta_i} - \frac{1}{c_\lambda + \lambda_j})\bigr) + \mu \Bigl( \frac{1}{c_\lambda} + \sum_{j=1}^m \frac{c_j}{\lambda_j} (\frac{1}{c_\lambda} - \frac{1}{\lambda_j + c_\lambda})\Bigr).
    \end{align*}
    
    To conclude moreover on the sign of the coefficient, it sufficies to note that as $c_\lambda > 0$:
    $$sign(\frac{1}{c_\lambda + \beta_i} - \frac{1}{c_\lambda + \lambda_j}) = sign(\lambda_j - \beta_i)$$
    It implies then that $\frac{1}{\lambda_j - \beta_i} (\frac{1}{c_\lambda + \beta_i} - \frac{1}{c_\lambda + \lambda_j}) > 0$ and by positivity of all $c_j$, all the terms involved in the sum are positive.
\end{proof}

\section{Proof of Section \ref{subsec:lob_sim}} \label{appendix:section_four_one}

\begin{proof}[Proof of Proposition \ref{prop:conditional_simulation}]
    \leavevmode \\
    We denote $q^\pi, L^\pi, C^\pi, N^\pi$ and $\mathcal{F}^\pi_\infty$ for the realization of the paths given a measure $\pi = (\pi^L, \pi^C, \pi^N)$.\\
    Conditional on $\mathcal{F}^\pi_\infty$, Theorem \ref{thm:conditionallaw} tells us that: 
    $$\mathcal{L}(\pi | \mathcal{F}^\pi_\infty) = \mathcal{L}(\pi^\sharp|\mathcal{F}^\pi_\infty), \, \pi^\sharp := (\pi^L_{|A_L^\pi}, \pi^C_{|A^\pi_C}, \pi^N_{|A^\pi_N}) \cup (\bar{\pi}^L_{|A^{\pi, c}_L}, \bar{\pi}^C_{|A^{\pi, c}_C}, \bar{\pi}^N_{|A^{\pi, c}_N})$$
    with $\pi^S_{|A^\pi_S} =^{\mathcal{L}} \sum_{n \geq 0} \ind_{(T_n^S, U_n^S \lambda^S(q^\pi_{T_n}))}$ with i.i.d. $U^S_n \sim \mathcal{U}([0,1])$, $\bar{\pi} = (\bar{\pi}^L, \bar{\pi}^C, \bar{\pi}^N)$ independent copy of $\pi$ and $A^\pi_x := \{ (s,z)|z \leq \lambda^x(q^\pi_s)\}$ for $x \in \{ L,C,N\}$.

    We want to describe the law of $\bar{q}^{\pi^\sharp}$, conditional on $\mathcal{F}^\pi_\infty$.

    Following the thinning description, one has $\bar{q}^{\pi^\sharp}_s := \bar{L}^{\pi^\sharp}_s - \bar{C}^{\pi^\sharp}_s - \bar{N}^{\pi^\sharp}_s + L^o_s$ and for $S \in \{L,C,N\}$:
    \begin{align*}
        \bar{S}^{\pi^\sharp}_t &= \int_0^t \int_{\R_+} \ind_{\{ z \leq \lambda^S(\bar{q}^{\pi^\sharp}_s) \}} \pi^{\sharp,S}(ds, dz) \\
        &= \int_0^t \int_{\R_+} \ind_{\{ z \leq \lambda^S(\bar{q}^{\pi^\sharp}_s)\} } \pi^S_{|A^\pi_S}(ds, dz) + \int_0^t \int_{\R_+} \ind_{\{ z \leq \lambda^S(\bar{q}^{\pi^\sharp}_s)\} } \bar{\pi}^S_{|A^{\pi,c}_S}(ds, dz) \\
        \bar{S}^{\tilde{\pi}}_t &= \int_0^t \int_0 \ind_{\{ z \leq \lambda^S(\bar{q}^{\pi^\sharp}_s)\} } \pi^S_{|A^\pi_S}(ds, dz) + \int_0^t \int_{\R_+} \ind_{\{\lambda^S(q^\pi_s) < z \leq \lambda^S(\bar{q}^{\pi^\sharp}_s)\} } \bar{\pi}^S(ds, dz)
    \end{align*}

    The event times for the simulation modifying the state $\bar{q}^{\pi^\sharp}$ come from:
    \begin{itemize}
        \item The first term of the equation: $\int_0^t \int_{\R_+} \ind_{\{ z \leq \lambda^S(\bar{q}^{\pi^\sharp}_s)\} } \pi^S_{|A^\pi_S}(ds, dz) = \sum_{T_n \leq t} \ind_{\{ U_n^S \lambda^S(q^\pi_{T_n}) \leq \lambda^S(\bar{q}^{\pi^\sharp}_{T_n})\} }$ for some $S \in \{ L,C,N\}$.
        \item The second term of the equation: $\int_0^t \int_{\R_+} \ind_{\{\lambda^S(q^\pi_s) < z \leq \lambda^S(\bar{q}^{\pi^\sharp}_s)\} } \bar{\pi}^S(ds, dz)$ for some $S \in \{L,C,N\}$.
        \item The metaorder execution: $L^o$.
    \end{itemize}

    Conditionally on the current state $\bar{q}^{\pi^\sharp}_t$ and $\mathcal{F}^\pi_\infty$, the next event of $\bar{q}^{\pi^\sharp}$ is the first event to occur out of the three sources mentioned.

    The terms $\sum_{T_n \leq t} \ind_{\{ U_n^S \lambda^S(q^\pi_{T_n}) \leq \lambda^S(\bar{q}^{\pi^\sharp}_{T_n})\} }$ give the times $t +\bar{\tau}^S$ for $S \in \{ L,C,N\}$. The terms $\int_0^t \int_{\R_+} \ind_{\{\lambda^S(q^\pi_s) < z \leq \lambda^S(\bar{q}^{\pi^\sharp}_s)\} } \bar{\pi}^S(ds, dz)$ for some $S \in \{L,C,N\}$ give the times $\tau^S$ for $S \in \{ L,C,N\}$ as the intensities are constant on between $t$ and the next event of $\bar{q}$, turning that equation in that of a homogeneous Poisson process for the next event time.
\end{proof}

\end{document}

